% !TeX root = ../main.tex
\section{Variational Derivation}
\label{sec:variational}

We derive the governing equations from Hamilton's extended principle, following
Bedford and Drumheller~\cite{drumheller1980thermomechanical}.  The formulation
is written for a compositional mixture containing a solid phase and an arbitrary
set $\mc M$ of mobile pore phases.  The symbols $l,k,w,n\in\mc M$ denote mobile
phases, while $\xi$ ranges over all phases.  Components within each phase are
indexed by $\alpha$.

\subsection{Variational principle and its terms}

The extended Hamilton's principle is
%
% AGENT-LOCK-BEGIN equation-lock: fixed through rendered equation (49); do not edit unless explicitly unlocked by user
\begin{equation*}
	\int_{t_1}^{t_2}\left[\delta(T - U) + \delta W + \sum_k \delta C_k\right] dt = 0,
\end{equation*}
% AGENT-LOCK-END
%
with kinetic energy, internal energy, virtual work, and four constraints:
%
% AGENT-LOCK-BEGIN equation-lock: fixed through rendered equation (49); do not edit unless explicitly unlocked by user
\begin{align}
	T        & = \sum_\xi \int_{\mc{B}}\frac{1}{2} \rho_\xi \mbf{v}_\xi \cdot \mbf{v}_\xi \dv,
	\label{eq:T_def}                                                                                                                                                                                                \\
	U        & = \sum_{\xi}\int_{\mc{B}} \rho_{\xi}\,\psi_{\xi}(\mbf{F}_{\xi}, \bar\rho_{\xi}, \eta^\al_{\xi})\dv
	+ \int_{\mc{B}} \phi\,\gamma\left(\left\{S_l\right\}_{l\in\mc M}, \theta\right) \dv,
	\label{eq:U_def}                                                                                                                                                                                                \\
	\delta W & = \sum_\xi \int_{\mc{B}} (\rho_\xi\mbf{g} + \mbf{b}_\xi)\cdot\delta\mbf{x}_\xi \,\dv
	-\sum_\xi\sum_\alpha\int_{\mc B}L_\xi^\alpha
	\frac{\delta\mc C_\xi^\alpha}{J_\xi}\,\dv
	- \sum_m\int_{\mc B}\mc A_{(m)}\,\delta\dot r_{(m)} \dv,
	\label{eq:W_def}                                                                                                                                                                                                \\
	C_1      & = \int_{\mc{B}} \lambda\left(\sum_\xi\phi_\xi - 1\right) \,\dv,
	\label{eq:C1_def}                                                                                                                                                                                               \\
	C_2      & = \sum_\xi \sum_\al \int_{\mc{B}_{\xi 0}} \tau^\al_\xi \left( J_\xi - \frac{\phi_{\xi 0} \eta^\al_{\xi 0} \bar\rho_{\xi 0} + \mc{C}_\xi^\alpha}{\phi_\xi \eta^\al_\xi \bar\rho_\xi} \right) \dV_\xi,
	\label{eq:C2_def}                                                                                                                                                                                               \\
	C_3      & = \sum_\xi \int_{\mc{B}} \pi_\xi \left( \sum_\alpha \eta_\xi^\alpha - 1 \right) \dv,
	\label{eq:C3_def}                                                                                                                                                                                               \\
	C_4      & = \sum_\xi\sum_\alpha\int_{\mc B}\mu_\xi^\alpha\left(\frac{\dot{\mc{C}}_\xi^\alpha}{J_\xi}-\sum_m\nu_{\xi (m)}^\alpha\dot r_{(m)}\right)\dv .
	\label{eq:C4_def}
\end{align}
% AGENT-LOCK-END
%
The coefficients $L_\xi^\alpha$ are component generalized forces conjugate to
adding component mass to phase $\xi$.  The virtual work includes gravity
$\mbf g$ and the net interphase interaction force $\mbf b_\xi$ applied to phase
$\xi$.  The affinities $\mc A_{(m)}$ are generalized mechanism forces conjugate
to the current progress-rate variations $\delta\dot r_{(m)}$.  Constraint $C_1$
enforces the volume-fraction constraint, $C_2$ enforces component mass balance,
$C_3$ enforces mass-fraction normalization within each phase, and $C_4$ ties
the mechanism progress rates to the phase-attached accumulated sources through
$\dot c_\xi^\alpha=\dot{\mc C}_\xi^\alpha/J_\xi$.  The multipliers
$\lambda$, $\tau_\xi^\alpha$, $\pi_\xi$, and $\mu_\xi^\alpha$ enforce these
constraints.
Chemical reactions and pure phase transformations are both represented by the
unified stoichiometric parameterization \eqref{eq:unified_stoich_current}.  The
fields $\mbf x_\xi$, $\eta_\xi^\alpha$, $\phi_\xi$, $\bar\rho_\xi$, and
$\mc C_\xi^\alpha$ are varied independently, while the mechanism rates
$\dot r_{(m)}$ are independent process-rate variables whose variations identify
the conjugate affinities.

\subsection{Extending Bedford and Drumheller identity 1.28}
\label{sec:extended_bd_identity}

Bedford and Drumheller~\cite[eq.~(1.28)]{drumheller1980thermomechanical}
give an identity for the variation of a constituent-attached spatial field
$\Gamma_\xi$ weighted by the current phase density $\rho_\xi$.  The same
transport argument extends componentwise to the compositional setting used here.

The accumulated source $\mc{C}_\xi^\alpha$ is attached to the phase-$\xi$
reference coordinates, while its current production rate is
$\dot c_\xi^\alpha=\dot{\mc{C}}_\xi^\alpha/J_\xi$.  Thus variations of source
terms must be pushed forward with the same phase map as the corresponding
component density.
%
For any phase field $\Gamma_\xi$, the Bedford--Drumheller identity extends componentwise as
%
% AGENT-LOCK-BEGIN equation-lock:  do not edit unless explicitly unlocked by user
\begin{align}
	\delta \int_{\mc{B}} \rho^\alpha_\xi \Gamma_\xi \, \dv
	 & = \delta \int_{\mc{B}_{\xi0}}
	\left(\rho^\alpha_{\xi0}+\mc{C}_\xi^\alpha\right)
	\Gamma_\xi \,\dVxi \notag        \\
	 & = \int_{\mc{B}_{\xi0}}
	\left[
		\delta\left(\rho^\alpha_{\xi0}+\mc{C}_\xi^\alpha\right)\Gamma_\xi
		+ J_\xi\rho^\alpha_\xi\delta\Gamma_\xi
	\right]\dVxi \notag              \\
	 & = \int_{\mc{B}_{\xi0}}
	\left[
		\delta\mc{C}_\xi^\alpha\Gamma_\xi
		+ J_\xi\rho^\alpha_\xi\delta\Gamma_\xi
	\right]\dVxi \notag              \\
	 & = \int_{\mc{B}}
	\left(
	\rho^\alpha_\xi\delta\Gamma_\xi
	+ \frac{\delta \mc{C}_\xi^\alpha}{J_\xi}\Gamma_\xi
	\right)\dv .
	\label{eq:bd128_mod2}
\end{align}
% AGENT-LOCK-END
%
Summing both sides of \eqref{eq:bd128_mod2} over $\alpha$ gives
%
% AGENT-LOCK-BEGIN equation-lock: do not edit unless explicitly unlocked by user
\begin{align}
	\delta \int_{\mc{B}} \rho_\xi \Gamma_\xi \, \dv
	 & = \int_{\mc{B}} \left(\rho_\xi \delta \Gamma_\xi + \sum_\alpha \frac{\delta \mc{C}^\al_\xi}{J_\xi} \Gamma_\xi  \right)  \dv
	\label{eq:bd128_mod}
\end{align}
% AGENT-LOCK-END
%
where $\rho_\xi = \sum_\alpha \rho^\alpha_\xi$.

% ============================================================
\subsection{Variation of the kinetic energy}
\label{sec:var_T_app}
% ============================================================

The variation of kinetic energy using \eqref{eq:bd128_mod} is then
%
% AGENT-LOCK-BEGIN equation-lock: fixed through rendered equation (49); do not edit unless explicitly unlocked by user
\begin{equation}
	\delta T = \sum_\xi \int_{\mc{B}} \rho_\xi\,\mbf{v}_\xi \cdot \delta\mbf{v}_\xi + \frac{1}{2} \mbf{v}_\xi \cdot \mbf{v}_\xi \sum_\al \frac{\delta \mc{C}^\al_\xi}{J_\xi} \dv.
\end{equation}
% AGENT-LOCK-END
%
Since $\delta\mbf{v}_\xi = \mathrm{d}(\delta\mbf{x}_\xi)/\mathrm{d}t$ and $\delta\mbf{x}_\xi\vert_{t_1} = \delta\mbf{x}_\xi\vert_{t_2} = \mbf{0}$, integration by parts in time gives
%
% AGENT-LOCK-BEGIN equation-lock: fixed through rendered equation (49); do not edit unless explicitly unlocked by user
\begin{equation}
	\int_{t_1}^{t_2} \delta T \dt
	= \int_{t_1}^{t_2} \sum_\xi \int_{\mc{B}} \left(-\rho_\xi\,\mbf{a}_\xi - \sum_\al \dot c^\alpha_\xi \mbf{v}_\xi \right) \cdot \delta\mbf{x}_\xi + \frac{1}{2} \mbf{v}_\xi \cdot \mbf{v}_\xi \sum_\al \frac{\delta \mc{C}_\xi^\alpha}{J_\xi}  \dv \dt
	\label{eq:deltaT_app}
\end{equation}
% AGENT-LOCK-END
%
\subsection{Variation of the phase free energies}
\label{sec:var_Uxi}
% ============================================================

For every phase, solid or mobile\footnote{The term \emph{mobile} denotes the
	nonsolid pore phases: aqueous, oleic, and gas.}, let $\psi_\xi$ denote the
specific Helmholtz free energy.  The corresponding contribution to the total
free energy is written
% AGENT-LOCK-BEGIN equation-lock: do not edit unless explicitly unlocked by user
\begin{align}
	U \supset U_\xi & = \sum_\xi \int_{\mc B}\rho_\xi\psi_\xi(\mbf F_\xi,\bar\rho_\xi,\eta_\xi^1, \eta_\xi^2, \dots, \eta_\xi^N)\,\dv \notag \\
	                & = \sum_\xi \int_{\mc B}\rho_\xi\psi_\xi(\mbf F_\xi,\bar\rho_\xi,\eta_\xi^\beta)\,\dv
	\label{eq:U_xi_uniform}
\end{align}
% AGENT-LOCK-END
Here $\eta_\xi^\beta$ denotes the full list of component mass fractions
$\{\eta_\xi^1,\ldots,\eta_\xi^N\}$; the index $\beta$ is not summed in this argument list.
The constitutive dependence on $\mbf F_\xi$ may be trivial for mobile phases,
but keeping it in the notation allows one calculation for all phases.  The local
free-energy variation is
%
% AGENT-LOCK-BEGIN equation-lock: fixed through rendered equation (49); do not edit unless explicitly unlocked by user
\begin{align}
	\delta\psi_\xi
	 & = \left.\p{\psi_\xi}{\mbf F_\xi}\right\vert_{\bar\rho_\xi,\eta_\xi^\beta}:\delta\mbf F_\xi
	+ \left.\p{\psi_\xi}{\bar\rho_\xi}\right\vert_{\mbf F_\xi,\eta_\xi^\beta}\delta\bar\rho_\xi
	+ \sum_\alpha
	\left.\p{\psi_\xi}{\eta_\xi^\alpha}\right\vert_{\mbf F_\xi,\bar\rho_\xi,\eta_\xi^{\beta\neq\alpha}}
	\delta\eta_\xi^\alpha .
	\label{eq:delta_psi_uniform}
\end{align}
% AGENT-LOCK-END
%
Using \eqref{eq:bd128_mod}, integrating the $\delta\mbf F_\xi$ term by parts,
and collecting independent variations gives
%
% AGENT-LOCK-BEGIN equation-lock: fixed through rendered equation (49); do not edit unless explicitly unlocked by user
\begin{align}
	\delta U_\xi
	 & = \sum_\xi\int_{\mc B}\Bigg[
		-\nabla_{\mbf x}\cdot\bs\sigma'_\xi\cdot\delta\mbf x_\xi
	+ \rho_\xi\left.\p{\psi_\xi}{\bar\rho_\xi}\right\vert_{\mbf F_\xi,\eta_\xi^\beta}\delta\bar\rho_\xi \notag \\
	 & \qquad
		+ \sum_\alpha \rho_\xi\left.\p{\psi_\xi}{\eta_\xi^\alpha}\right\vert_{\mbf F_\xi,\bar\rho_\xi,\eta_\xi^{\beta\neq\alpha}}\delta\eta_\xi^\alpha
		+ \sum_\alpha \psi_\xi\frac{\delta \mc{C}_\xi^\alpha}{J_\xi}
		\Bigg]\dv + \text{boundary terms}.
	\label{eq:deltaU_uniform}
\end{align}
% AGENT-LOCK-END
%
The corresponding Cauchy effective stress is
%
% AGENT-LOCK-BEGIN equation-lock: fixed through rendered equation (49); do not edit unless explicitly unlocked by user
\begin{align}
	\bs\sigma'_\xi
	\equiv \rho_\xi\left.\p{\psi_\xi}{\mbf F_\xi}\right\vert_{\bar\rho_\xi,\eta_\xi^\beta}\mbf F_\xi^T .
	\label{eq:var_cauchy_def}
\end{align}
% AGENT-LOCK-END
%
For mobile phases this term vanishes when $\psi_\xi$ has no dependence on
$\mbf F_\xi$.



% ============================================================
\subsection{Variation of the interfacial energies}
\label{sec:var_gamma}
% ============================================================

Let $\mc I$ denote the set of oriented mobile interfaces.  Each
$(w,n)\in\mc I$ identifies a wetting phase $w$ and a non-wetting phase $n$.
The total mobile volume fraction and the mobile saturations are
% AGENT-LOCK-BEGIN equation-lock: do not edit unless explicitly unlocked by user
\begin{equation}
	\phi = \sum_{l\in\mc M}\phi_l,
	\qquad
	S_l = \frac{\phi_l}{\phi},
	\qquad
	\sum_{l\in\mc M}S_l=1 .
	\label{eq:mobile_saturation_definitions}
\end{equation}
% AGENT-LOCK-END
The saturations are used here only as ratios of the primitive volume-fraction
fields; the variations are still taken with respect to the primitive
$\phi_l$.  The notation $\left\{S_l\right\}_{l\in\mc M}$ denotes the mobile
saturation state subject to the constraint in
\eqref{eq:mobile_saturation_definitions}.  Thus the energy may be written in
terms of any independent set of mobile saturations, with the remaining
saturation determined by the constraint; the derivation below keeps the
constraint explicit instead of choosing a reference mobile phase.  For thermal
reservoir simulation the interfacial energy is allowed to depend on
temperature, consistent with the standard capillary-pressure scaling by
interfacial tension in reservoir engineering \cite{leverett1941capillary}.  The
effective mobile interfacial energy density is a single constitutive potential
of this constrained saturation state and temperature, and it is not assumed to
decompose into additive pair contributions:
% AGENT-LOCK-BEGIN equation-lock: do not edit unless explicitly unlocked by user
\begin{equation}
	\gamma =
	\gamma\left(\left\{S_l\right\}_{l\in\mc M},\theta\right).
	\label{eq:interfacial_energy_density}
\end{equation}
% AGENT-LOCK-END
and the corresponding contribution to the internal energy is
%
% AGENT-LOCK-BEGIN equation-lock: do not edit unless explicitly unlocked by user
\begin{equation}
	U \supset U_\gamma =
	\int_{\mc{B}}
	\phi\,\gamma\left(\left\{S_l\right\}_{l\in\mc M},\theta\right)
	\dv .
	\label{eq:interfacial_internal_energy}
\end{equation}
% AGENT-LOCK-END

Before varying the interfacial energy, record the transported variation of the
phase-attached volume fractions.  The Bedford--Drumheller spatial identity
(cf.\ \cite{drumheller1980thermomechanical} equation (1.25)) gives, for each
mobile phase field $\phi_l$,
%
% AGENT-LOCK-BEGIN equation-lock: do not edit unless explicitly unlocked by user
\begin{equation}
	\left.\delta\phi_\xi\right\vert_{\mbf{x}}
	= \delta\phi_\xi-\nabla_{\mbf x}\phi_\xi\cdot\delta\mbf x_\xi.
	\label{eq:transported_volume_fraction_variation}
\end{equation}
% AGENT-LOCK-END
%
Using
% AGENT-LOCK-BEGIN equation-lock: do not edit unless explicitly unlocked by user
\begin{equation}
	\delta\phi=\sum_{l\in\mc M}\delta\phi_l,
	\qquad
	\delta S_l=\frac{\delta\phi_l-S_l\delta\phi}{\phi},
	\label{eq:saturation_variations}
\end{equation}
% AGENT-LOCK-END
the fixed-point contribution and the transport correction from
\eqref{eq:transported_volume_fraction_variation} are collected in one
calculation:
%
% AGENT-LOCK-BEGIN equation-lock: fixed through rendered equation (49); do not edit unless explicitly unlocked by user
\begin{align}
	\delta U_\gamma
	 & = \int_{\mc B}\left[
	\gamma\,\delta\phi
	+\sum_{l\in\mc M}
	\left.\p{\gamma}{S_l}\right\vert_{\theta,S_{k\neq l}}
	\left(\delta\phi_l-S_l\delta\phi\right)
	\right]\dv \notag                                          \\
	 & \quad
	-\int_{\mc B}\sum_{l\in\mc M}
	\gamma_l\nabla_{\mbf x}\phi_l\cdot\delta\mbf x_l\dv \notag \\
	 & = \int_{\mc B}\sum_{l\in\mc M}\gamma_l\delta\phi_l\dv
	-\int_{\mc B}\sum_{l\in\mc M}\gamma_l\nabla_{\mbf x}\phi_l\cdot\delta\mbf x_l\dv .
	\label{eq:deltaU_fg}
\end{align}
% AGENT-LOCK-END
%
The phase interfacial potential appearing in \eqref{eq:deltaU_fg} is
%
% AGENT-LOCK-BEGIN equation-lock: fixed through rendered equation (49); do not edit unless explicitly unlocked by user
\begin{equation}
	\gamma_l
	=
	\gamma
	+\left.\p{\gamma}{S_l}\right\vert_{\theta,S_{k\neq l}}
	-\sum_{k\in\mc M}S_k
	\left.\p{\gamma}{S_k}\right\vert_{\theta,S_{j\neq k}},
	\label{eq:primitive_gamma_derivatives}
\end{equation}
% AGENT-LOCK-END
%
where the saturation derivatives are evaluated on the saturation simplex.  The
oriented wetting/non-wetting pairs specify capillary-pressure differences, not a
preferred set of independent saturation coordinates.  The wetting/non-wetting
interfacial difference associated with an oriented pair is
%
% AGENT-LOCK-BEGIN equation-lock: fixed through rendered equation (49); do not edit unless explicitly unlocked by user
\begin{equation}
	p_{c,nw}\equiv\gamma_n-\gamma_w,
	\qquad (w,n)\in\mc I .
	\label{eq:pc_def}
\end{equation}
% AGENT-LOCK-END
%
The form \eqref{eq:deltaU_fg} keeps the gradient terms generated by the scalar
transport of the phase-attached volume fractions.  These terms combine with the
volume-fraction and component-storage multiplier gradients in the phase momentum
balances after the volume-fraction Euler--Lagrange equations are used.

% ============================================================
\subsection{Variation of constraint \texorpdfstring{$C_1$}{C1}: volume-fraction}
\label{sec:var_C1_app}
% ============================================================

The constraint $C_1 = \int_{\mc B} \lambda(\sum_\xi\phi_\xi - 1)\,\dv$ is
written in the current configuration, so the natural variation is Eulerian,
taken at fixed spatial point $\mbf x$.  Since each $\phi_\xi$ is a
constituent-attached spatial field, \eqref{eq:transported_volume_fraction_variation}
expresses the variation in terms of the spatial variation $\delta\phi_\xi$ and
the motion variation $\delta\mbf x_\xi$:
%
% AGENT-LOCK-BEGIN equation-lock: fixed through rendered equation (49); do not edit unless explicitly unlocked by user
\begin{align}
	\delta C_1
	 & = \int_{\mc{B}} \delta\lambda \left(\sum_\xi\phi_\xi - 1\right) \dv
	+ \sum_\xi \int_{\mc{B}} \lambda\left(\delta\phi_\xi - \nabla_{\mbf{x}}\phi_\xi\cdot\delta\mbf{x}_\xi\right) \dv.
	\label{eq:deltaC1_app}
\end{align}
% AGENT-LOCK-END
%
The first term enforces the volume-fraction constraint
$\sum_\xi\phi_\xi = 1$.  The remaining terms contribute $\lambda$ to each
$\delta\phi_\xi$ equation and $-\lambda\nabla_{\mbf x}\phi_\xi$ to each
momentum variation $\delta\mbf x_\xi$.

% ============================================================
\subsection{Variation of constraint \texorpdfstring{$C_2$}{C2}: component material mass balance}
\label{sec:var_C2_app}
% ============================================================

The component material mass balance constraint is
%
% AGENT-LOCK-BEGIN equation-lock: fixed through rendered equation (49); do not edit unless explicitly unlocked by user
\begin{equation}
	C_2 = \sum_\xi\sum_\alpha\int_{\mc{B}_{\xi0}}\tau_\xi^\alpha \left(J_\xi - \frac{\phi_{\xi0}\bar\rho_{\xi0} \eta^\al_{\xi 0} + \mc{C}_\xi^\alpha}{\phi_\xi\bar\rho_\xi \eta_\xi^\alpha} \right) \,dV_\xi.
	\label{eq:C2_component_measure}
\end{equation}
% AGENT-LOCK-END
%
Hence
%
% AGENT-LOCK-BEGIN equation-lock: fixed through rendered equation (49); do not edit unless explicitly unlocked by user
\begin{align}
	\delta C_2
	 & = \sum_\xi\sum_\alpha \int_{\mc{B}_{\xi0}}
	\delta\tau_\xi^\alpha
	\left(J_\xi - \frac{\phi_{\xi0}\bar\rho_{\xi0}\eta_{\xi0}^\alpha + \mc{C}_\xi^\alpha}{\phi_\xi\bar\rho_\xi\eta_\xi^\alpha}\right)
	\,{\rm d}V_\xi
	\notag                                        \\
	 & \quad
	+ \sum_\xi\sum_\alpha \int_{\mc{B}_{\xi0}}
	\tau_\xi^\alpha \delta J_\xi \, {\rm d}V_\xi
	\notag                                        \\
	 & \quad
	+ \sum_\xi\sum_\alpha \int_{\mc{B}_{\xi0}}
	\tau_\xi^\alpha\,\frac{\phi_{\xi0}\bar\rho_{\xi0}\eta_{\xi0}^\alpha + \mc{C}_\xi^\alpha}{\phi_\xi^2\bar\rho_\xi \eta_\xi^\alpha}\,\delta\phi_\xi \, {\rm d}V_\xi
	\notag                                        \\
	 & \quad
	+ \sum_\xi\sum_\alpha \int_{\mc{B}_{\xi0}}
	\tau_\xi^\alpha\,\frac{\phi_{\xi0}\bar\rho_{\xi0}\eta_{\xi0}^\alpha + \mc{C}_\xi^\alpha}{\phi_\xi \bar\rho_\xi^2 \eta^\al_\xi}\,\delta\bar\rho_\xi \, {\rm d}V_\xi
	\notag                                        \\
	 & \quad
	+ \sum_\xi\sum_\alpha \int_{\mc{B}_{\xi0}}
	\tau_\xi^\alpha\,\frac{\phi_{\xi0}\bar\rho_{\xi0}\eta_{\xi0}^\alpha + \mc{C}_\xi^\alpha}{\phi_\xi \bar\rho_\xi (\eta^\al_\xi)^2}\,\delta\eta_\xi^\alpha \, {\rm d}V_\xi
	\notag                                        \\
	 & \quad
	- \sum_\xi\sum_\alpha \int_{\mc{B}_{\xi0}}
	\frac{\tau_\xi^\alpha}{\phi_\xi\bar\rho_\xi \eta_\xi^\alpha}\,\delta \mc{C}_\xi^\alpha \, {\rm d}V_\xi.
	\notag                                        \\
	 & = \sum_\xi\sum_\alpha \int_{\mc{B}_{\xi0}}
	\delta\tau_\xi^\alpha
	\left(J_\xi - \frac{\phi_{\xi0}\bar\rho_{\xi0}\eta_{\xi0}^\alpha + \mc{C}_\xi^\alpha}{\phi_\xi\bar\rho_\xi\eta_\xi^\alpha}\right)
	\,{\rm d}V_\xi
	\notag                                        \\
	 & \quad
	-\sum_\xi\sum_\alpha \int_{\mc{B}}
	\nabla_{\mbf{x}}\tau_\xi^\alpha \cdot \delta\mbf{x}_\xi \, \dv
	\notag                                        \\
	 & \quad
	+ \sum_\xi\sum_\alpha\int_{\mc{B}}
	\left[
		\frac{\tau_\xi^\alpha}{\phi_\xi}\delta\phi_\xi
		+ \frac{\tau_\xi^\alpha}{\bar\rho_\xi}\delta\bar\rho_\xi
		+ \frac{\tau_\xi^\alpha}{\eta_\xi^\alpha}\delta\eta_\xi^\alpha
		- \frac{\tau_\xi^\alpha}{\phi_\xi\bar\rho_\xi\eta_\xi^\alpha}\frac{\delta \mc{C}_\xi^\alpha}{J_\xi}
		\right]\dv
	+ \text{b.t.}
	\label{eq:deltaC2_lines}
\end{align}
% AGENT-LOCK-END
%
Here the second equality uses $\dv = J_\xi {\rm d}V_\xi$ and
$\phi_{\xi0}\bar\rho_{\xi0}\eta_{\xi0}^\alpha+\mc{C}_\xi^\alpha
	=J_\xi\phi_\xi\bar\rho_\xi\eta_\xi^\alpha$.

% ============================================================
\subsection{Variation of constraint \texorpdfstring{$C_3$}{C3}: mass partition within phase}
\label{sec:var_C3_app}
% ============================================================

%
% AGENT-LOCK-BEGIN equation-lock: fixed through rendered equation (49); do not edit unless explicitly unlocked by user
\begin{align}
	\delta C_3
	 & = \sum_\xi \int_{\mc{B}} \delta\pi_\xi \left(\sum_\al \eta_\xi^\al - 1\right) \dv
	+ \sum_\xi \sum_\al \int_{\mc{B}} \pi_\xi \left(\delta\eta^\al_\xi - \nabla_{\mbf{x}}\eta_\xi^\al\cdot\delta\mbf{x}_\xi\right) \dv, \notag \\
	 & = \sum_\xi \int_{\mc{B}} \delta\pi_\xi \left(\sum_\al \eta_\xi^\al - 1\right) \dv
	+ \sum_\xi \sum_\al \int_{\mc{B}} \pi_\xi \delta\eta^\al_\xi \dv.
	\label{eq:deltaC3_app}
\end{align}
% AGENT-LOCK-END
%
The last term in the first line vanishes by recognizing $\sum_\alpha \nabla_{\mbf{x}} \eta_\xi^\alpha = 0$ by \eqref{eq:mass_fraction_constraint}.
% ============================================================
% ============================================================
\subsection{Stoichiometric conversion variations}
\label{sec:stoich_conversion_variations}
% ============================================================
Conversion sources are restricted by the current-rate kinematic constraint
\eqref{eq:C4_def}.
The variation of $\dot c_\xi^\alpha=\dot{\mc{C}}_\xi^\alpha/J_\xi$ must be taken before
using the stoichiometric projection, because the denominator depends on the phase
map.  Repeating the transport calculation used for phase-attached rates gives
%
% AGENT-LOCK-BEGIN equation-lock: fixed through rendered equation (49); do not edit unless explicitly unlocked by user
\begin{align}
	\int_{t_1}^{t_2}\delta C_4\,\dt
	 & =\int_{t_1}^{t_2}\int_{\mc B}\sum_\xi\sum_\alpha
	\delta\mu_\xi^\alpha\left(\dot c_\xi^\alpha-
	\sum_m\nu_{\xi (m)}^\alpha\dot r_{(m)}\right)\dv\dt \notag \\
	 & \quad+\int_{t_1}^{t_2}\int_{\mc B}\sum_\xi\sum_\alpha
	\left[-\frac{D_\xi\mu_\xi^\alpha}{Dt}\frac{\delta \mc{C}_\xi^\alpha}{J_\xi}
		+\dot c_\xi^\alpha\nabla_{\mbf x}\mu_\xi^\alpha\cdot\delta\mbf x_\xi\right]
	\dv\dt \notag                                              \\
	 & \quad-\int_{t_1}^{t_2}\int_{\mc B}\sum_m
	\left(\sum_\xi\sum_\alpha\mu_\xi^\alpha\nu_{\xi (m)}^\alpha\right)
	\delta\dot r_{(m)}\,\dv\dt+\text{b.t.}
	\label{eq:C4_variation}
\end{align}
% AGENT-LOCK-END
%
% ============================================================
\subsection{Collected variational principle}
\label{sec:collected_var}
% ============================================================

Collecting the coefficients of the variations $\delta\mbf{x}_\xi$, $\delta\phi_\xi$, $\delta\bar\rho_\xi$, $\delta\eta^\alpha_\xi$, together with the independent process-rate variations $\delta\dot r_{(m)}$, gives the following collection before substituting any Euler--Lagrange equation back into another variation.
%
% AGENT-LOCK-BEGIN equation-lock:  do not edit unless explicitly unlocked by user
\begin{align}
	0 = & \sum_{\xi} \int_{\mc{B}} \Bigg[
	-\rho_{\xi}\mbf{a}_{\xi}
	+ \nabla_{\mbf{x}}\cdot\bs\sigma'_{\xi} + \rho_{\xi}\mbf{g} + \mbf{b}_{\xi}
	-\lambda\nabla_{\mbf{x}}\phi_\xi
	-\sum_\alpha\nabla_{\mbf{x}}\tau_\xi^\alpha
	+ \sum_\alpha \dot c_{\xi}^{\alpha}\left(\nabla_{\mbf x}\mu_\xi^\alpha-\mbf{v}_{\xi}\right)
	\Bigg]\cdot\delta\mbf{x}_{\xi} \,\dv \notag                                                                                                              \\
	    & + \int_{\mc B}\Big[
		\sum_{l\in\mc M}
		\gamma_l\nabla_{\mbf x}\phi_l\cdot\delta\mbf x_l
	\Big]\dv \notag                                                                                                                                          \\
	    & + \sum_\xi \int_{\mc{B}} \Bigg[
		\lambda + \sum_\alpha \frac{\tau_{\xi}^\alpha}{\phi_{\xi}}
	\Bigg] \delta\phi_{\xi} \,\dv \notag                                                                                                                     \\
	    & - \int_{\mc B}\sum_{l\in\mc M}\gamma_l\delta\phi_l\dv \notag                                                                                       \\
	    & + \sum_{\xi}\int_{\mc{B}} \Bigg[
		-\rho_{\xi}\left.\p{\psi_{\xi}}{\bar\rho_{\xi}}\right\vert_{\mbf{F}_{\xi},\eta_\xi^\beta}
		+ \sum_\alpha\frac{\tau_{\xi}^\alpha}{\bar\rho_{\xi}}
	\Bigg] \delta\bar\rho_{\xi} \,\dv \notag                                                                                                                 \\
	    & + \sum_{\xi}\sum_\alpha \int_{\mc{B}} \Bigg[
		\frac{\tau_{\xi}^\alpha}{\eta_\xi^\alpha} + \pi_{\xi}
		- \rho_{\xi}\left.\p{\psi_{\xi}}{\eta^\alpha_\xi}\right\vert_{\mbf F_\xi,\bar\rho_{\xi},\eta^{\beta\neq\alpha}_\xi}
	\Bigg] \delta\eta^\alpha_\xi \,\dv \notag                                                                                                                \\
	    & + \sum_\xi\sum_\alpha\int_{\mc B}\left[
		\tfrac12\mbf v_\xi\cdot\mbf v_\xi-\psi_\xi-L_\xi^\alpha
		-\frac{D_\xi\mu_\xi^\alpha}{Dt}
		-\frac{\tau_\xi^\alpha}{\phi_\xi\bar\rho_\xi\eta_\xi^\alpha}
	\right]\frac{\delta \mc{C}_\xi^\alpha}{J_\xi}\,\dv \notag                                                                                                \\
	    & - \sum_m\int_{\mc B}\left(\mc A_{(m)}+\sum_\xi\sum_\alpha\mu_\xi^\alpha\nu_{\xi (m)}^\alpha\right)\delta\dot r_{(m)}\,\dv + \text{boundary terms}.
	\label{eq:collected_full}
\end{align}
% AGENT-LOCK-END
%
The multiplier variations are omitted here because they return the constraint equations~\eqref{eq:C1_def}--\eqref{eq:C4_def}. Since all variations are considered arbitrary, each bracketed integrand must vanish pointwise, yielding the Euler--Lagrange equations.

\subsection{Euler--Lagrange equations}

The Euler--Lagrange equations generated by arbitrary $\delta\phi_\xi$ give,
for phases without an explicit capillary contribution,
%
% AGENT-LOCK-BEGIN equation-lock: fixed through rendered equation (49); do not edit unless explicitly unlocked by user
\begin{equation}
	\phi_\xi \lambda = - \sum_\alpha \tau^\al_\xi .
	\label{eq:el_lambda}
\end{equation}
% AGENT-LOCK-END
%
For the mobile phases the corresponding equations contain the primitive
capillary potentials discussed below.

Substituting \eqref{eq:el_lambda} into the Euler-Lagrange equations associated with arbitrary $\delta \mbf{x}_\xi$ and simplifying gives
%
% AGENT-LOCK-BEGIN equation-lock: fixed through rendered equation (49); do not edit unless explicitly unlocked by user
\begin{equation}
	\rho_{\xi}\mbf{a}_{\xi} =
	\nabla_{\mbf{x}}\cdot\bs\sigma'_{\xi} + \rho_{\xi}\mbf{g} + \mbf{b}_{\xi}
	+ \phi_{\xi} \nabla_{\mbf{x}}\lambda
	+ \sum_\alpha \dot c_{\xi}^{\alpha}\left(\nabla_{\mbf x}\mu_\xi^\alpha-\mbf{v}_{\xi}\right),
	\label{eq:el_momentum_uniform}
\end{equation}
% AGENT-LOCK-END
%
with the mobile-phase capillary-force additions displayed separately below.

The Euler--Lagrange equation associated with arbitrary $\delta\bar\rho_\xi$ gives
%
% AGENT-LOCK-BEGIN equation-lock: fixed through rendered equation (49); do not edit unless explicitly unlocked by user
\begin{equation}
	\sum_\al \frac{\tau_{\xi}^\al}{\phi_{\xi}}
	= \boxed{\bar\rho^2_{\xi}\left.\p{\psi_{\xi}}{\bar\rho_{\xi}}\right\vert_{\mbf F_\xi,\eta_\xi^\beta}
	\equiv \bar{p}_{\xi}} .
	\label{eq:phase_pressure_all}
\end{equation}
% AGENT-LOCK-END
%
and defines the thermodynamic phase pressures $\bar{p}_\xi$.

Let $l\in\mc M$ denote a mobile phase.
The coefficient of $\delta\phi_l$ in \eqref{eq:collected_full} includes the
primitive interfacial potential $\gamma_l$, so the mobile-phase volume-fraction
relation is
%
% AGENT-LOCK-BEGIN equation-lock: do not edit unless explicitly unlocked by user
\begin{equation}
	\lambda+\sum_\alpha\frac{\tau_l^\alpha}{\phi_l}-\gamma_l=0.
	\label{eq:mobile_volume_fraction_el_relation}
\end{equation}
% AGENT-LOCK-END
%
Using \eqref{eq:phase_pressure_all} gives
%
% AGENT-LOCK-BEGIN equation-lock: do not edit unless explicitly unlocked by user
\begin{equation}
	\lambda+\bar p_l-\gamma_l=0,
	\qquad\hbox{or}\qquad
	\bar p_l=\gamma_l-\lambda .
	\label{eq:mobile_pressure_primitive_potential_relation}
\end{equation}
% AGENT-LOCK-END
%
Consequently pressure differences between mobile phases are exactly the
differences of the primitive interfacial potentials,
%
% AGENT-LOCK-BEGIN equation-lock: do not edit unless explicitly unlocked by user
\begin{equation}
	\bar p_l-\bar p_k=\gamma_l-\gamma_k .
	\label{eq:mobile_pressure_difference_primitive_potential}
\end{equation}
% AGENT-LOCK-END
%
Using \eqref{eq:pc_def}, the wetting/non-wetting capillary pressure relation is
recovered:
%
% AGENT-LOCK-BEGIN equation-lock: do not edit unless explicitly unlocked by user
\begin{equation}
	\boxed{p_{c,nw}=\bar p_n-\bar p_w},
	\qquad (w,n)\in\mc I .
	\label{eq:wetting_nonwetting_capillary_pressure_recovery}
\end{equation}
% AGENT-LOCK-END
%
For two mobile phases with $\mc M=\{w,n\}$ and one oriented pair $(w,n)$, this
reduces to the usual two-phase capillary pressure.  If the interfacial energy is
parameterized by the wetting saturation $S_w$, then
$p_{c,nw}=-\partial\gamma/\partial S_w$.

The same relations identify the equivalent pore pressure associated with the
volume constraint.  Define
%
% AGENT-LOCK-BEGIN equation-lock: do not edit unless explicitly unlocked by user
\begin{equation}
	\bar p_E\equiv-\lambda .
	\label{eq:equivalent_pore_pressure_definition}
\end{equation}
% AGENT-LOCK-END
%
Multiplying $\bar p_l=\gamma_l+\bar p_E$ by the mobile saturations and summing
over mobile phases gives
%
% AGENT-LOCK-BEGIN equation-lock: do not edit unless explicitly unlocked by user
\begin{equation}
	\sum_l S_l\bar p_l
	=
	\bar p_E+\sum_l S_l\gamma_l .
	\label{eq:saturation_weighted_pressure_sum}
\end{equation}
% AGENT-LOCK-END
%
The primitive potentials satisfy
%
% AGENT-LOCK-BEGIN equation-lock: do not edit unless explicitly unlocked by user
\begin{equation}
	\sum_l S_l\gamma_l
	=
	\gamma .
	\label{eq:saturation_weighted_primitive_potential_sum}
\end{equation}
% AGENT-LOCK-END
%
Therefore
%
% AGENT-LOCK-BEGIN equation-lock: do not edit unless explicitly unlocked by user
\begin{equation}
	\boxed{\bar p_E=\sum_l S_l\bar p_l-\gamma} .
	\label{eq:equivalent_pore_pressure_capillary_relation}
\end{equation}
% AGENT-LOCK-END
%
For a phase with $N$ tracked components, the Euler--Lagrange equations
associated with arbitrary $\delta\eta_{\xi}^\alpha$ give
%
% AGENT-LOCK-BEGIN equation-lock: fixed through rendered equation (49); do not edit unless explicitly unlocked by user
\begin{equation}
	0 = \frac{\tau_{\xi}^\alpha}{\eta_{\xi}^\alpha} + \pi_{\xi}
	- \rho_{\xi}\left.\p{\psi_{\xi}}{\eta_{\xi}^\alpha}\right\vert_{\mbf F_\xi,\bar\rho_{\xi},\eta_{\xi}^{\beta\neq\alpha}},
	\qquad \alpha=1,\ldots,N,
	\label{eq:el_eta_all}
\end{equation}
% AGENT-LOCK-END
%
while variation with respect to $\pi_{\xi}$ returns the composition normalization
constraint
%
% AGENT-LOCK-BEGIN equation-lock: fixed through rendered equation (49); do not edit unless explicitly unlocked by user
\begin{equation}
	\sum_{\alpha=1}^N \eta_{\xi}^\alpha = 1.
	\label{eq:el_eta_constraint}
\end{equation}
% AGENT-LOCK-END
%
Thus $\pi_{\xi}$ is the phase-level multiplier that enforces the redundant
composition coordinate.  Equation~\eqref{eq:el_eta_all} may be read
componentwise as
%
% AGENT-LOCK-BEGIN equation-lock: fixed through rendered equation (49); do not edit unless explicitly unlocked by user
\begin{equation}
	\rho_{\xi}\left.\p{\psi_{\xi}}{\eta_{\xi}^\alpha}\right\vert_{\mbf F_\xi,\bar\rho_{\xi},\eta_{\xi}^{\beta\neq\alpha}}
	- \frac{\tau_{\xi}^\alpha}{\eta_{\xi}^\alpha}
	= \pi_{\xi},
	\label{eq:pi_componentwise}
\end{equation}
% AGENT-LOCK-END
%
% AGENT-LOCK-BEGIN equation-lock: fixed through rendered equation (49); do not edit unless explicitly unlocked by user
\begin{align}
	\rho_{\xi}\left(
	\left.\p{\psi_{\xi}}{\eta_{\xi}^\alpha}\right\vert_{\mbf F_\xi,\bar\rho_{\xi},\eta_{\xi}^{\kappa\neq\alpha}}
	- \left.\p{\psi_{\xi}}{\eta_{\xi}^\beta}\right\vert_{\mbf F_\xi,\bar\rho_{\xi},\eta_{\xi}^{\kappa\neq\beta}}
	\right)
	 & = \frac{\tau_{\xi}^\alpha}{\eta_{\xi}^\alpha}
	- \frac{\tau_{\xi}^\beta}{\eta_{\xi}^\beta}.
	\label{eq:pi_difference}
\end{align}
% AGENT-LOCK-END
%
Here $\kappa$ denotes the collection of composition components held fixed in the
partial derivative; it is not summed.

Choosing component $N$ as the reference and comparing against
$\beta=1,\ldots,N-1$ gives a system of $N-1$ equations,
%
% AGENT-LOCK-BEGIN equation-lock: fixed through rendered equation (49); do not edit unless explicitly unlocked by user
\begin{align}
	\rho_{\xi}\left(
	\left.\p{\psi_{\xi}}{\eta_{\xi}^\beta}\right\vert_{\mbf F_\xi,\bar\rho_{\xi},\eta_{\xi}^{\kappa\neq \beta}}
	- \left.\p{\psi_{\xi}}{\eta_{\xi}^N}\right\vert_{\mbf F_\xi,\bar\rho_{\xi},\eta_{\xi}^{\kappa\neq N}}
	\right)
	 & = \frac{\tau_{\xi}^\beta}{\eta_{\xi}^\beta}
	- \frac{\tau_{\xi}^N}{\eta_{\xi}^N},
	\qquad \beta=1,\ldots,N-1.
	\label{eq:pi_difference_system}
\end{align}
% AGENT-LOCK-END
%
The equations associated with the conversion variables are
%
% AGENT-LOCK-BEGIN equation-lock: fixed through rendered equation (49); do not edit unless explicitly unlocked by user
\begin{align}
	\frac{D_\xi\mu_\xi^\alpha}{Dt}
	            & = \frac{1}{2}\mbf v_\xi\cdot\mbf v_\xi
	-\psi_\xi
	-L_\xi^\alpha
	-\frac{\tau_\xi^\alpha}{\phi_\xi\bar\rho_\xi\eta_\xi^\alpha},
	\label{eq:el_conversion_component}                                     \\
	\mc A_{(m)} & =-\sum_\xi\sum_\alpha\mu_\xi^\alpha\nu_{\xi (m)}^\alpha,
	\qquad m=1,\ldots,N_{(m)}.
	\label{eq:el_conversion_affinity}
\end{align}
% AGENT-LOCK-END
%
Equation~\eqref{eq:el_conversion_affinity} is not an additional closure relation
for the state variables.  It defines the mechanism affinity as the stoichiometric
projection of the phase-attached conversion potentials, but only after the
kinematic push-forward $\dot c_\xi^\alpha=\dot{\mc{C}}_\xi^\alpha/J_\xi$ has been
varied.  The closure enters only when one prescribes either an equilibrium
condition $\mc A_{(m)}=0$ or a kinetic relation
$\dot r_{(m)}=\mc K_{(m)}(\mc A_{(m)},\text{state})$.  Thus, for kinetically limited
mechanisms such as slow fluid--solid reactions, \eqref{eq:el_conversion_affinity}
supplies the driving force rather than an algebraic equilibrium equation.  For an
instantaneous phase transformation, $\dot r_{(m)}$ may be a
measure in time.  In all cases, the phase momentum equations retain the inertial
production term and the $\nabla\mu_\xi^\alpha$ force generated by adding or
removing mass in a phase moving with velocity $\mbf v_\xi$.  For a
single-component solid, $\eta_{\xi_s}^\alpha=1$.

At this variational stage $L_\xi^\alpha$ remains an explicit constitutive
conversion-work potential.  An admissible choice of $L_\xi^\alpha$ is identified
by the Coleman--Noll argument in
Section~\ref{sec:constitutive_models_entropy_inequality}.

\subsection{Mobile-phase momentum equations}
\label{sec:mobile_phase_momentum}

For any mobile phase $l\in\mc M$, the scalar-transport capillary contribution
gives the intermediate form
%
% AGENT-LOCK-BEGIN equation-lock: do not edit unless explicitly unlocked by user
\begin{align}
	\rho_l\mbf a_l
	 & = \nabla_{\mbf x}\cdot\bs\sigma'_l + \rho_l\mbf g + \mbf b_l
	-\lambda\nabla_{\mbf x}\phi_l
	-\sum_\alpha\nabla_{\mbf x}\tau_l^\alpha
	+\gamma_l\nabla_{\mbf x}\phi_l
	+ \sum_\alpha \dot c_l^{\alpha}\left(\nabla_{\mbf x}\mu_l^\alpha-\mbf{v}_l\right).
	\label{eq:el_mom_l_start}
\end{align}
% AGENT-LOCK-END
%
Using $\sum_\alpha\tau_l^\alpha=\phi_l\bar p_l$ and
$\lambda+\bar p_l-\gamma_l=0$, the multiplier and interfacial-gradient terms
collapse as
%
% AGENT-LOCK-BEGIN equation-lock: do not edit unless explicitly unlocked by user
\begin{equation}
	-\lambda\nabla_{\mbf x}\phi_l
	-\nabla_{\mbf x}\left(\phi_l\bar p_l\right)
	+\gamma_l\nabla_{\mbf x}\phi_l
	=
	-\phi_l\nabla_{\mbf x}\bar p_l .
	\label{eq:mobile_capillary_pressure_gradient_collapse}
\end{equation}
% AGENT-LOCK-END
%
The standard phase pressure-gradient form is therefore
%
% AGENT-LOCK-BEGIN equation-lock: do not edit unless explicitly unlocked by user
\begin{empheq}[box=\fbox]{equation}
	\rho_l\mbf{a}_l
	= \rho_l\mbf{g} + \mbf{b}_l
	- \phi_l\nabla_{\mbf{x}}\bar p_l
	+ \sum_\alpha \dot c_l^{\alpha}\left(\nabla_{\mbf x}\mu_l^\alpha-\mbf{v}_l\right),
	\qquad l\in\mc M .
	\label{eq:el_mom_l_simplified}
\end{empheq}
% AGENT-LOCK-END
%

\section{Constitutive Models and the Entropy Inequality}
\label{sec:multicomponent_solid_mixture_rules}
\label{sec:constitutive_models_entropy_inequality}

The variational derivation above admits multicomponent phases.  For the
constitutive specialization used in the Coleman--Noll argument, we now restrict
the solid portion of the mixture to a set $\mc S$ of mineral solid phases.  Each
solid phase carries one active component label, retained in the superscript for
stoichiometric bookkeeping, and satisfies
\(\eta_s^s=1\).  Thus the component partial density and phase density coincide,
\(\rho_s^s=\rho_s\), while reactions still identify the precipitating or
dissolving solid through the phase/component pair \((s,s)\).  All solid phases
are attached to a common skeleton trajectory
\(\mbf x=\bs\chi(\mbf X,t)\), with deformation gradient \(\mbf F\) and Jacobian
\(J\).

For the present fully compressible solid model we still need the distension
kinematics before it is meaningful to define an intrinsic stress.  The common
skeleton has macroscopic deformation \(\mbf F\), but the stress carried by each
solid phase is conjugate to that phase's true, or undistended, deformation.  We
therefore introduce a solid-phase distension gradient tensor \(\mbf A_s\) and
write
%
% AGENT-LOCK-BEGIN equation-lock: do not edit unless explicitly unlocked by user
\begin{equation}
	\mbf F=\mbf A_s\bar{\mbf F}_s,
	\qquad
	J=a_s\bar J_s,
	\qquad
	a_s=\det\mbf A_s,
	\qquad
	\bar J_s=\det\bar{\mbf F}_s,
	\qquad
	s\in\mc S .
	\label{eq:solid_distension_decomposition}
\end{equation}
% AGENT-LOCK-END
%
Here $\bar{\mbf F}_s$ is the true deformation seen by solid phase \(s\), while
$\mbf A_s$ accounts for the phase-specific pore-space allocation, packing, and
stress-free volume accommodation between the common skeleton motion and the true
solid deformation.

Following Drumheller's treatment of reacting immiscible mixtures with
elastic--plastic true deformation and distension~\cite{drumheller2000}, we
split the full distension tensor into elastic, plastic, and reversible
thermal--compositional pore-architecture factors,
% AGENT-LOCK-BEGIN equation-lock: do not edit unless explicitly unlocked by user
\begin{equation}
	\mbf A_s=\mbf A_s^e\mbf A_s^p
	\mbf A_s^\theta,
	\qquad
	s\in\mc S .
	\label{eq:solid_tensor_distension_split}
\end{equation}
% AGENT-LOCK-END
The tensor $\mbf A_s^\theta=\mbf A_s^\theta(\theta)$ is a reversible stress-free
pore-architecture, or pore-volume, factor describing the temperature-dependent
change of pore volume, packing, fabric, and solid allocation associated with
solid phase \(s\).  No explicit functional form is prescribed at this level.
Thermo-poroelastic treatments of heated saturated rocks make precisely this
distinction by assigning a thermal expansion coefficient to the pore volume
separately from the solid grains, so that thermal pressurization depends on the
mismatch among fluid, pore-volume, and grain expansions~\cite{palciauskas1982thermally,ghabezloo2009stress,ghabezloo2011micromechanics}.
Here that pore-volume part is carried by $\mbf A_s^\theta$, while
$\bar{\mbf F}_s^\theta$ below carries the stress-free intrinsic thermal
dilation of the true solid phase.  The tensor $\mbf A_s^e$ stores reversible
distension energy, while $\mbf A_s^p$ is an internal variable that records
permanent pore-architecture, fabric, or compaction change.

The distension tensor is the kinematic device used to separate deformation of
the solid material from changes in pore-space allocation, packing, and
reaction-induced filling.  The evolution of $a_s$ follows from mass balance
together with the decomposition $J=a_s\bar J_s$.  For each one-component solid
phase, the material balance gives
%
% AGENT-LOCK-BEGIN equation-lock: do not edit unless explicitly unlocked by user
\begin{equation}
	J\phi_s\bar\rho_s^s
	=\phi_{s0}\bar\rho_{s0}^s+\mc{C}_s^s,
	\qquad
	\eta_s^s=1,
	\qquad
	s\in\mc S .
	\label{eq:solid_material_mass_distension}
\end{equation}
% AGENT-LOCK-END
%
Using $J=a_s\bar J_s$ and imposing true local conservation of mass for the
solid material,
%
% AGENT-LOCK-BEGIN equation-lock: do not edit unless explicitly unlocked by user
\begin{equation}
	\bar J_s\bar\rho_s^s=\bar\rho_{s0}^s,
	\qquad
	s\in\mc S .
	\label{eq:solid_true_mass_conservation}
\end{equation}
% AGENT-LOCK-END
%
transforms \eqref{eq:solid_material_mass_distension} into
%
% AGENT-LOCK-BEGIN equation-lock: do not edit unless explicitly unlocked by user
\begin{equation}
	a_s\phi_s\bar\rho_{s0}^s
	=\phi_{s0}\bar\rho_{s0}^s+\mc{C}_s^s,
	\qquad
	s\in\mc S .
	\label{eq:solid_distension_mass_relation}
\end{equation}
% AGENT-LOCK-END
%
Thus the solid volume fraction is tied directly to the motion only when there is
no mass exchange.  The rate split follows from the Jacobian product in
\eqref{eq:solid_distension_decomposition}: taking the natural logarithm gives
$\ln J=\ln a_s+\ln\bar J_s$, and then taking the material time derivative gives
%
% AGENT-LOCK-BEGIN equation-lock: do not edit unless explicitly unlocked by user
\begin{empheq}[box=\fbox]{equation}
	\frac{\dot a_s}{a_s}
	=\frac{\dot J}{J}-\frac{\dot{\bar J}_s}{\bar J_s} .
	\label{eq:solid_distension_jacobian_rate_split}
\end{empheq}
% AGENT-LOCK-END
%
Since $\dot J=J\nabla_{\mbf x}\cdot\mbf v_s$ and
\eqref{eq:solid_true_mass_conservation} implies
${\dot{\bar J}_s}/{\bar J_s}=-{\dot{\bar\rho}_s^s}/{\bar\rho_s^s}$, the distension
evolution law is
%
% AGENT-LOCK-BEGIN equation-lock: do not edit unless explicitly unlocked by user
\begin{equation}
	\frac{\dot a_s}{a_s}
	=\nabla_{\mbf x}\cdot\mbf v_s
	+\frac{\dot{\bar\rho}_s^s}{\bar\rho_s^s},
	\qquad
	s\in\mc S .
	\label{eq:solid_distension_evolution}
\end{equation}
% AGENT-LOCK-END
%

For non-isothermal elastic--plastic problems each solid phase has its own
stress-free true deformation, plastic true deformation, and elastic true
deformation, while remaining attached to the common skeleton trajectory:
%
% AGENT-LOCK-BEGIN equation-lock: do not edit unless explicitly unlocked by user
\begin{equation}
	\bar{\mbf F}_s
	=\bar{\mbf F}_s^e\bar{\mbf F}_s^p\bar{\mbf F}_s^\theta,
	\qquad
	s\in\mc S .
	\label{eq:true_elastic_plastic_thermal_split}
\end{equation}
% AGENT-LOCK-END
%
Here $\bar{\mbf F}_s^\theta$ maps to the thermal stress-free true-solid
configuration of solid phase \(s\), $\bar{\mbf F}_s^p$ is the plastic
true-deformation internal variable whose rate is closed by a constitutive flow
rule, and
$\bar{\mbf F}_s^e
	=\bar{\mbf F}_s(\bar{\mbf F}_s^\theta)^{-1}(\bar{\mbf F}_s^p)^{-1}$
is the dependent elastic part that enters the intrinsic stresses.
Following Drumheller~\cite{drumheller2000}, the plastic
true-deformation flow is isochoric,
% AGENT-LOCK-BEGIN equation-lock: do not edit unless explicitly unlocked by user
\begin{equation*}
	\det\bar{\mbf F}_s^p=1,
	\qquad
	\operatorname{tr}\left[
		\dot{\bar{\mbf F}}_s^p(\bar{\mbf F}_s^p)^{-1}
		\right]=0,
	\qquad
	s\in\mc S .
\end{equation*}
% AGENT-LOCK-END
A yield surface, cap surface, critical-state law, or reaction-assisted
compaction law can then be introduced as a constitutive specialization of the
residual dissipation.  The stress-free true-solid thermal/compositional factor
is prescribed generically as
$\bar{\mbf F}_s^\theta=\bar{\mbf F}_s^\theta(\theta)$.  Explicit isotropic or
anisotropic thermal expansion laws are constitutive specializations.

This thermal split does not break the energy-based structure.  It is a
constitutive reparametrization of the free energy: instead of depending on
$\bar{\mbf F}$ directly, the elastic part of each solid phase free energy
depends on $\bar{\mbf F}_s^e$ and $\mbf A_s^e$, while $\bar{\mbf F}_s^p$ and
$\mbf A_s^p$ enter as internal variables through the kinematic rates.  The
superscript label is retained for reaction and density bookkeeping, but the
solid Helmholtz free energy is written with phase-level notation in the
thermodynamic derivation:
%
% AGENT-LOCK-BEGIN equation-lock: do not edit unless explicitly unlocked by user
\begin{equation}
	\psi_s
	=\psi_s(\bar{\mbf F}_s^e,\mbf A_s^e,\bar\rho_s^s,\theta),
	\qquad
	\psi_s\equiv\psi_s^s,
	\qquad
	\eta_s^s=1,
	\qquad
	s\in\mc S .
	\label{eq:solid_multicomponent_free_energy}
\end{equation}
% AGENT-LOCK-END
%
Temperature dependence of the free energy supplies entropy and heat capacity in
the usual way, while the dependence of $\bar{\mbf F}_s^\theta$ on $\theta$ and
the dependence of $\mbf A_s^\theta$ on $\theta$ encode stress-free intrinsic thermal
expansion and stress-free distension for each solid phase.  The thermodynamic
restrictions that make these constitutive choices admissible are obtained next
from the entropy inequality.

\section{Coleman--Noll procedure for the compositional thermo-reactive mixture}
\label{sec:coleman_noll_compositional}

This appendix records the thermodynamic argument in a form that mirrors the
kinematic split used in the body of the paper.  The point is not to postulate the
solid mixture rules, but to show how they appear after the component Helmholtz
energies are differentiated subject to the density, composition, and dilation
constraints.

The Coleman--Noll calculation is applied on the reduced constitutive state space.
It is not a second derivation of the constrained equations of motion.  Constraints
that define the admissible motions and fields, such as saturation, storage, and
stoichiometric source constraints, are enforced in the variational formulation
and do not introduce additional Coleman--Noll coefficients unless the constrained
variable is an argument of a constitutive free energy.  Following the standard
Liu multiplier exploitation of constrained entropy inequalities~\cite{liu1972method},
we may add identically vanishing constraint-rate terms before collecting
coefficients.  Composition enters here because the free energies depend explicitly
on $\eta_\xi^\alpha$; its rate is therefore restricted to the tangent space of
the mass-fraction constraint,
\begin{equation}
	\sum_\alpha\eta_\xi^\alpha=1,
	\qquad
	\sum_\alpha\dot\eta_\xi^\alpha=0 .
	\label{eq:CN_composition_tangent_constraint}
\end{equation}

For each phase/component pair let
\begin{equation}
	e_\xi^\alpha=\psi_\xi^\alpha+\theta s_\xi^\alpha
	\label{eq:CN_legendre}
\end{equation}
be the Legendre relation.  We write each phase free energy as a mass-specific
phase potential
\begin{equation}
	\psi_\xi=\sum_\alpha\eta_\xi^\alpha\psi_\xi^\alpha,
	\qquad
	\rho_\xi\psi_\xi=\sum_\alpha\rho_\xi^\alpha\psi_\xi^\alpha .
	\label{eq:CN_phase_free_energy_sum}
\end{equation}
For the solid phase we also introduce the component volume fraction within the
solid phase, denoted $\phi_s^\alpha$.  It is not an independent thermodynamic
state variable in the Coleman--Noll derivatives; it is a derived volume-weight
used later to express skeleton mixture rules:
\begin{equation}
	\rho_s^\alpha
	=\phi_s\bar\rho_s\eta_s^\alpha
	=\phi_s\phi_s^\alpha\bar\rho_s^\alpha,
	\qquad
	\sum_\alpha\eta_s^\alpha=1,
	\qquad
	\sum_\alpha\phi_s^\alpha=1 .
	\label{eq:CN_solid_component_volume_fraction}
\end{equation}
Thus
\begin{equation}
	\bar\rho_s=\sum_\alpha\phi_s^\alpha\bar\rho_s^\alpha,
	\qquad
	\eta_s^\alpha=\frac{\phi_s^\alpha\bar\rho_s^\alpha}{\bar\rho_s},
	\qquad
	\phi_s^\alpha=\frac{\eta_s^\alpha\bar\rho_s}{\bar\rho_s^\alpha} .
	\label{eq:CN_eta_phi_volume_relation}
\end{equation}
The primitive solid component Helmholtz free energies are taken to have the same
density and composition state variables as the mobile phases,
\begin{equation}
	\mathcal{S} = \{\bar{\mbf F}_s^M, a_s, \bar\rho_s,
	\theta,\eta_s^\beta \}
\end{equation}
%


\begin{align}
	\psi_s^\alpha
	 & =\psi_s^\alpha(\bar{\mbf F}_s^M, a_s, \bar\rho_s,
	\theta,\eta_s^\beta),
	\qquad
	\bar{\mbf F}_s^M=\bar{\mbf F}_s(\bar{\mbf F}_s^\theta)^{-1},
	\label{eq:CN_component_free_energy_arguments}
\end{align}
where
$\bar{\mbf F}_s^\theta=\bar{\mbf F}_s^\theta(\theta,\phi_s^\beta)$ is a
prescribed stress-free thermal/compositional dilation, not an independent
thermodynamic state.  Since the Coleman--Noll state variables remain
$(\bar\rho_s,\eta_s^\beta)$, derivatives of $\bar{\mbf F}_s^\theta$ with respect
to composition are total derivatives through the derived component volume
fractions.  Define the state sets
\begin{equation}
	\mc X\in\{\mc S,\mc F,\mc T_s\},
	\qquad
	\mc T_s=\{\theta,\phi_s^\beta\}.
	\label{eq:CN_fixed_eta_state_set}
\end{equation}
When a derivative is written for a generic phase, let
\begin{equation}
	\mc X_\xi =
	\begin{cases}
		\mc S, & \xi=s,    \\
		\mc F, & \xi=\ell,
	\end{cases}
	\label{eq:CN_phase_state_set}
\end{equation}
We use the convention that
\begin{equation}
	\left.\p{\Gamma_\xi}{z}\heldfixed{\mc X}
	\label{eq:CN_held_fixed_convention}
\end{equation}
denotes differentiation with respect to $z$ while holding fixed all other
members of the state set $\mc X$.
We write
\begin{equation}
	\left.\p{\bar{\mbf F}_s^\theta}{\eta_s^\beta}\heldfixed{\mc S}
	=\sum_\gamma
	\left.\p{\bar{\mbf F}_s^\theta}{\phi_s^\gamma}\heldfixed{\mc T_s}
	\left.\p{\phi_s^\gamma}{\eta_s^\beta}\heldfixed{\mc S} .
	\label{eq:CN_Ftheta_eta_chain}
\end{equation}
For the mobile fluid phases, we postulate the Coleman--Noll state variables
\begin{equation}
	\mc F
	=\left\{\bar\rho_\ell,\eta_\ell^\beta,\theta\right\},
	\qquad
	\psi_\ell^\alpha=\psi_\ell^\alpha(\mc F) .
	\label{eq:CN_fluid_state_variables}
\end{equation}
Thus the fluid-phase free-energy rates entering the entropy inequality are
\begin{align}
	\sum_\ell\sum_\alpha\rho_\ell^\alpha\dot\psi_\ell^\alpha
	 & =\sum_\ell\rho_\ell\sum_\alpha\eta_\ell^\alpha\left(
	\left.\p{\psi_\ell^\alpha}{\bar\rho_\ell}\heldfixed{\mc F}\dot{\bar\rho}_\ell
	+\sum_\beta\left.\p{\psi_\ell^\alpha}{\eta_\ell^\beta}\heldfixed{\mc F}\dot\eta_\ell^\beta
	+\left.\p{\psi_\ell^\alpha}{\theta}\heldfixed{\mc F}\dot\theta
	\right).
	\label{eq:CN_fluid_free_energy_rate}
\end{align}
The dependent quantities
\begin{equation}
	\rho_\xi^\alpha=\phi_\xi\bar\rho_\xi\eta_\xi^\alpha,
	\qquad
	\sum_\alpha\eta_\xi^\alpha=1,
	\qquad
	\bar J_s=\det\bar{\mbf F}_s=\frac{J_s}{a_s}
	\label{eq:CN_kinematic_dependencies}
\end{equation}

For the solid, component additivity gives
\begin{align}
	\sum_\alpha\rho_s^\alpha\dot{\psi}_s^\alpha
	 & =\rho_s\sum_\alpha\eta_s^\alpha\dot\psi_s^\alpha \notag \\
	 & =\rho_s\sum_\alpha\eta_s^\alpha \bigg(
	\left.\p{\psi_s^\alpha}{\bar{\mbf{F}}^M_s}\heldfixed{\mc S}
	:\dot{\bar{\mbf{F}}}^M_s
	+\left.\p{\psi_s^\alpha}{a_s}\heldfixed{\mc S}\dot{a}_s
	\notag                                                     \\
	 & \qquad
	+\left.\p{\psi_s^\alpha}{\bar\rho_s}\heldfixed{\mc S}\dot{\bar\rho}_s
	+\sum_\beta\left.\p{\psi_s^\alpha}{\eta_s^\beta}\heldfixed{\mc S}\dot\eta_s^\beta
	\notag                                                     \\
	 & \qquad
	+\left.\p{\psi_s^\alpha}{\theta}\heldfixed{\mc S}\dot\theta
	\bigg).
	\label{eq:CN_total_helmholtz_density}
\end{align}
The Coleman--Noll argument treats $\dot{\bar{\mbf F}}_s$ as the independent
solid deformation rate.  The mechanical deformation gradient is a dependent
constitutive argument, so substituting the prescribed thermal/compositional
stress-free dilation rate gives
\begin{align}
	\dot{\bar{\mbf F}}_s^M
	= & \; \dot{\bar{\mbf F}}_s(\bar{\mbf F}_s^\theta)^{-1}
	-\bar{\mbf F}_s^M
	\left.\frac{\partial\bar{\mbf F}_s^\theta}{\partial\theta}\heldfixed{\mc T_s}
	(\bar{\mbf F}_s^\theta)^{-1}\dot\theta \notag           \\
	  & \;-
	\sum_\beta
	\bar{\mbf F}_s^M
	\left.\p{\bar{\mbf F}_s^\theta}{\eta_s^\beta}\heldfixed{\mc S}
	(\bar{\mbf F}_s^\theta)^{-1}\dot\eta_s^\beta .
	\label{eq:CN_FM_rate_chain_rule}
\end{align}
It is useful to define the tensor conjugate to the stress-free dilation,
\begin{equation}
	\mbf H_s^\alpha
	= (\bar{\mbf F}_s^M)^T
	\left.\p{\psi_s^\alpha}{\bar{\mbf F}_s^M}\heldfixed{\mc S}
	(\bar{\mbf F}_s^\theta)^{-T} .
	\label{eq:CN_H_def}
\end{equation}
The mobile-phase Helmholtz energies are phase EOS potentials using the
fluid-state variables in \eqref{eq:CN_fluid_state_variables}.  We assign no
elastic stress power to the mobile phases, so
\begin{equation}
	\bs\sigma_\ell'=\mbf 0 .
	\label{eq:CN_mobile_stress_free}
\end{equation}
\begin{align}
	\sum_\alpha\left(\rho_\xi^\alpha\dot e_\xi^\alpha
	+\dot c_\xi^\alpha e_\xi^\alpha\right)
	= & \; \bs\sigma_\xi':\mbf L_\xi
	+\phi_\xi\bar p_\xi\frac{\dot{\bar\rho}_\xi}{\bar\rho_\xi}
	+\sum_\alpha\frac{\tau_\xi^\alpha}{\eta_\xi^\alpha}\dot\eta_\xi^\alpha
	-\nabla_{\mbf x}\cdot\mbf q_\xi
	+\rho_\xi r_\xi+\rho_\xi e_\xi^{\rm int}.
	\label{eq:CN_energy_balance}
\end{align}
The second constraint-power term is the compositional analog of the pressure
power.  The multipliers $\tau_\xi^\alpha$ enforce component material storage in
\eqref{eq:C2_def}; when the free energy depends on the normalized composition,
the admissible composition rate does work against the normalized storage
multiplier $\tau_\xi^\alpha/\eta_\xi^\alpha$.
For the solid phase, the mechanical power in \eqref{eq:CN_energy_balance} must
be decomposed with the same kinematic split used in the constitutive law.  In the
isotropic dilation specialization,
\begin{equation}
	\mbf F_s=a_s^{1/3}\bar{\mbf F}_s,
	\label{eq:CN_F_split_for_power}
\end{equation}
so
\begin{equation}
	\mbf L_s
	=\dot{\mbf F}_s\mbf F_s^{-1}
	=\frac{1}{3}\frac{\dot a_s}{a_s}\mbf I
	+\dot{\bar{\mbf F}}_s\bar{\mbf F}_s^{-1} .
	\label{eq:CN_L_split_for_power}
\end{equation}
Therefore
\begin{equation}
	\bs\sigma_s':\mbf L_s
	=\left[\bs\sigma_s'\bar{\mbf F}_s^{-T}\right]
	:\dot{\bar{\mbf F}}_s
	+\frac{1}{3a_s}\operatorname{tr}(\bs\sigma_s')\dot a_s .
	\label{eq:CN_stress_power_split}
\end{equation}
The entropy inequality is
\begin{equation}
	0\le
	\sum_\xi\left[\sum_\alpha\left(
		\rho_\xi^\alpha\dot s_\xi^\alpha
		+\dot c_\xi^\alpha s_\xi^\alpha\right)
		+\nabla_{\mbf x}\cdot\left(\frac{\mbf q_\xi}{\theta}\right)
		-\frac{\rho_\xi r_\xi}{\theta}\right].
	\label{eq:CN_entropy_inequality}
\end{equation}
Using \eqref{eq:CN_energy_balance} to eliminate $\rho_\xi r_\xi$ from
\eqref{eq:CN_entropy_inequality}, and then using
$e_\xi^\alpha=\psi_\xi^\alpha+\theta s_\xi^\alpha$, gives the reduced entropy
inequality.  At this stage we also add the identically vanishing rate of the
composition-normalization constraint,
$\pi_\xi\sum_\alpha\dot\eta_\xi^\alpha$, so that the Coleman--Noll coefficient
collection is performed on the tangent space
\eqref{eq:CN_composition_tangent_constraint}.
\begin{align}
	0\le
	\sum_\xi \bigg[ &
		\frac{1}{\theta}\Big(
		\bs\sigma_\xi':\mbf L_\xi
		+\phi_\xi\bar p_\xi\frac{\dot{\bar\rho}_\xi}{\bar\rho_\xi}
		+\sum_\alpha\frac{\tau_\xi^\alpha}{\eta_\xi^\alpha}\dot\eta_\xi^\alpha
		+\pi_\xi\sum_\alpha\dot\eta_\xi^\alpha
		\notag                                                \\
		& \qquad
		-\sum_\alpha\rho_\xi^\alpha
		\left(\dot\psi_\xi^\alpha+s_\xi^\alpha\dot\theta\right)
		-\sum_\alpha\dot c_\xi^\alpha\psi_\xi^\alpha
		+\rho_\xi e_\xi^{\rm int}\Big)
		-\frac{\mbf q_\xi\cdot\nabla_{\mbf x}\theta}{\theta^2}
		\bigg].
	\label{eq:CN_reduced_entropy_inequality}
\end{align}
Substituting \eqref{eq:CN_total_helmholtz_density},
\eqref{eq:CN_FM_rate_chain_rule}, \eqref{eq:CN_stress_power_split}, and
\eqref{eq:CN_fluid_free_energy_rate} gives the collected solid--fluid
contribution.
\begin{align}
	0\le\; &
	\frac{1}{\theta}\Bigg\{
	\left[
	\bs\sigma_s'\bar{\mbf F}_s^{-T}
	-\rho_s\sum_\alpha\eta_s^\alpha
	\left.\p{\psi_s^\alpha}{\bar{\mbf F}_s^M}\heldfixed{\mc S}
	(\bar{\mbf F}_s^\theta)^{-T}
	\right]:\dot{\bar{\mbf F}}_s
	\notag                                               \\
	       & +\left[
		\frac{1}{3a_s}\operatorname{tr}(\bs\sigma_s')
		-\rho_s\sum_\alpha\eta_s^\alpha
		\left.\p{\psi_s^\alpha}{a_s}\heldfixed{\mc S}
		\right]\dot a_s
	\notag                                               \\
	       & +\sum_\xi \left[
		\frac{\phi_\xi \bar p_\xi}{\bar\rho_\xi}
		-\rho_\xi\sum_\alpha\eta_\xi^\alpha
		\left.\p{\psi_\xi^\alpha}{\bar\rho_\xi}\heldfixed{\mc X_\xi}
		\right]\dot{\bar\rho}_\xi
		\notag                                               \\
	       & -\sum_\xi\rho_\xi\sum_\alpha\eta_\xi^\alpha
		\left[
			s_\xi^\alpha+\left.\p{\psi_\xi^\alpha}{\theta}\heldfixed{\mc X_\xi}
			\right]\dot\theta
	+\rho_s\sum_\alpha\eta_s^\alpha
	\mbf H_s^\alpha:
	\left.\p{\bar{\mbf F}_s^\theta}{\theta}\heldfixed{\mc T_s}\dot\theta
	\notag                                               \\
	       & -\sum_\xi\rho_\xi\sum_\beta
		\left.\p{\psi_\xi}{\eta_\xi^\beta}\heldfixed{\mc X_\xi}
		\dot\eta_\xi^\beta
		+\rho_s\sum_\beta\left[
		\sum_\alpha\eta_s^\alpha\mbf H_s^\alpha:
		\left.\p{\bar{\mbf F}_s^\theta}{\eta_s^\beta}\heldfixed{\mc S}
		\right]\dot\eta_s^\beta
		\notag                                               \\
	       & +\sum_\xi\sum_\beta\frac{\tau_\xi^\beta}{\eta_\xi^\beta}
		\dot\eta_\xi^\beta
		+\sum_\xi \pi_\xi\sum_\beta\dot\eta_\xi^\beta
		\notag                                               \\
	       & -\sum_m\left(
		\sum_\xi\sum_\alpha\nu_{\xi m}^\alpha\psi_\xi^\alpha
		\right)\dot r_m
		\notag                                               \\
	       & +\sum_\xi\rho_\xi e_\xi^{\rm int}
	-\sum_\xi\frac{\mbf q_\xi\cdot\nabla_{\mbf x}\theta}{\theta}
	\Bigg\} .
	\label{eq:CN_solid_collected_entropy}
\end{align}
The phase derivative with respect to composition is the derivative of the phase
Helmholtz potential in \eqref{eq:CN_phase_free_energy_sum}.  For mobile phases no
elastic stress-free dilation term appears, so the composition coefficient gives
the same stationarity condition obtained from Hamilton's principle in
\eqref{eq:el_eta_all}:
\begin{equation}
	0=
	\frac{\tau_\ell^\beta}{\eta_\ell^\beta}
	+\pi_\ell
	-\rho_\ell
	\left.\p{\psi_\ell}{\eta_\ell^\beta}\heldfixed{\mc F},
	\qquad
	\beta=1,\ldots,N .
	\label{eq:CN_composition_multiplier_restriction}
\end{equation}
Equivalently,
\begin{equation}
	\rho_\ell
	\left.\p{\psi_\ell}{\eta_\ell^\beta}\heldfixed{\mc F}
	-\frac{\tau_\ell^\beta}{\eta_\ell^\beta}
	=\pi_\ell .
	\label{eq:CN_composition_componentwise}
\end{equation}
For the solid phase, the same coefficient collection contains the explicit
stress-free compositional dilation correction from
\eqref{eq:CN_FM_rate_chain_rule}:
\begin{equation}
	0=
	\frac{\tau_s^\beta}{\eta_s^\beta}
	+\pi_s
	-\rho_s
	\left.\p{\psi_s}{\eta_s^\beta}\heldfixed{\mc S}
	+\rho_s\sum_\alpha\eta_s^\alpha\mbf H_s^\alpha:
	\left.\p{\bar{\mbf F}_s^\theta}{\eta_s^\beta}\heldfixed{\mc S},
	\qquad
	\beta=1,\ldots,N .
	\label{eq:CN_solid_composition_multiplier_restriction}
\end{equation}
Eliminating $\pi_\ell$ gives the mobile-phase projected restriction on the
tangent space,
\begin{equation}
	\rho_\ell\left(
	\left.\p{\psi_\ell}{\eta_\ell^\beta}\heldfixed{\mc F}
	-\left.\p{\psi_\ell}{\eta_\ell^N}\heldfixed{\mc F}
	\right)
	=
	\frac{\tau_\ell^\beta}{\eta_\ell^\beta}
	-\frac{\tau_\ell^N}{\eta_\ell^N},
	\qquad
	\beta=1,\ldots,N-1 .
	\label{eq:CN_composition_projected_restriction}
\end{equation}
For the solid, eliminating $\pi_s$ gives the corresponding projected restriction
with the difference of the dilation corrections retained:
\begin{align}
	 & \rho_s\left(
	\left.\p{\psi_s}{\eta_s^\beta}\heldfixed{\mc S}
	-\left.\p{\psi_s}{\eta_s^N}\heldfixed{\mc S}
	\right)
	-\rho_s\sum_\alpha\eta_s^\alpha\mbf H_s^\alpha:
	\left(
	\left.\p{\bar{\mbf F}_s^\theta}{\eta_s^\beta}\heldfixed{\mc S}
	-\left.\p{\bar{\mbf F}_s^\theta}{\eta_s^N}\heldfixed{\mc S}
	\right)
	\notag                                                     \\
	 & \qquad =
	\frac{\tau_s^\beta}{\eta_s^\beta}
	-\frac{\tau_s^N}{\eta_s^N},
	\qquad
	\beta=1,\ldots,N-1 .
	\label{eq:CN_solid_composition_projected_restriction}
\end{align}
The deformation coefficient gives
\begin{align}
	\bs\sigma_s'
	 & =\rho_s\sum_\alpha\eta_s^\alpha
	\left.\p{\psi_s^\alpha}{\bar{\mbf F}_s^M}\heldfixed{\mc S}
	(\bar{\mbf F}_s^M)^T
	=\phi_s\sum_\alpha\phi_s^\alpha\bar\rho_s^\alpha
	\left.\p{\psi_s^\alpha}{\bar{\mbf F}_s^M}\heldfixed{\mc S}
	(\bar{\mbf F}_s^M)^T .
	\label{eq:CN_stress_definition}
\end{align}
The intrinsic skeleton stress may therefore be written as the component volume
fraction average
\begin{equation}
	\bar{\bs\sigma}_s'
	=\frac{1}{\phi_s}\bs\sigma_s'
	=\sum_\alpha\phi_s^\alpha\bar{\bs\sigma}_s^{\prime\alpha},
	\qquad
	\bar{\bs\sigma}_s^{\prime\alpha}
	=\bar\rho_s^\alpha
	\left.\p{\psi_s^\alpha}{\bar{\mbf F}_s^M}\heldfixed{\mc S}
	(\bar{\mbf F}_s^M)^T .
	\label{eq:CN_intrinsic_component_stress}
\end{equation}
Equivalently, if $\bar\rho_{s0}^\alpha\psi_s^\alpha$ is the energy per
undistended intrinsic reference volume of component $\alpha$, then
\begin{equation}
	\bar{\mbf P}_s^{\prime\alpha}
	=\bar\rho_{s0}^\alpha
	\left.\p{\psi_s^\alpha}{\bar{\mbf F}_s^M}\heldfixed{\mc S}
	=\bar J_s \bar{\bs\sigma}_s^{\prime\alpha}
	(\bar{\mbf F}_s^M)^{-T}
	\label{eq:CN_intrinsic_pk_stress}
\end{equation}
is the undistended intrinsic first Piola--Kirchhoff stress.  The distension and
pressure coefficients give
\begin{align}
	\frac{1}{3}\operatorname{tr}(\bs\sigma_s')
	 & = \rho_s a_s \sum_\alpha\eta_s^\alpha
	\left.\p{\psi_s^\alpha}{a_s}\heldfixed{\mc S},
	\label{eq:CN_distension_force_definition}             \\
	\bar p_\xi
	 & =\left(\bar\rho_\xi\right)^2\sum_\alpha\eta_\xi^\alpha
	\left.\p{\psi_\xi^\alpha}{\bar\rho_\xi}\heldfixed{\mc X_\xi} .
	\label{eq:CN_phase_pressure_definition}
\end{align}
The entropy coefficient gives
\begin{equation}
	\sum_\alpha\eta_s^\alpha s_s^\alpha
	=-\sum_\alpha\eta_s^\alpha
	\left.\p{\psi_s^\alpha}{\theta}\heldfixed{\mc S}
	+\sum_\alpha\eta_s^\alpha
	\mbf H_s^\alpha:
	\left.\p{\bar{\mbf F}_s^\theta}{\theta}\heldfixed{\mc T_s} .
	\label{eq:CN_entropy_definitions}
\end{equation}
The solid composition restrictions above are the place where stress-free
compositional dilation modifies the usual composition coefficient.  This form
gives the desired skeleton mixture rules.  If the Legendre-transformed
intrinsic law for each component has the small-strain volumetric form
\begin{equation}
	\bar\sigma_{m,s}^{\prime\prime\alpha}
	=K_\alpha^\star\varepsilon_v-b_\alpha\hat p,
	\label{eq:CN_component_skeleton_linearization}
\end{equation}
then $\bar\sigma_{m,s}^{\prime\prime}=\sum_\alpha\phi_s^\alpha
	\bar\sigma_{m,s}^{\prime\prime\alpha}$ implies
\begin{equation}
	K^\star=\sum_\alpha\phi_s^\alpha K_\alpha^\star,
	\qquad
	b=\sum_\alpha\phi_s^\alpha b_\alpha .
	\label{eq:CN_skeleton_mixture_rules}
\end{equation}
The mixture-volume stress contribution remains
$\bs\sigma_s^{\prime\prime}=\phi_s\bar{\bs\sigma}_s^{\prime\prime}$.

Finally, using $\dot c_\xi^\alpha=\sum_m\nu_{\xi m}^\alpha\dot r_m$, the
conversion virtual power from the variational derivation contributes the
multiplier part of the affinity projection,
\begin{equation}
	\sum_m\mc A_m\dot r_m
	+\sum_\xi\sum_\alpha\mu_\xi^\alpha
	\sum_m\nu_{\xi m}^\alpha\dot r_m
	=\sum_m\left(\mc A_m+
	\sum_\xi\sum_\alpha\mu_\xi^\alpha\nu_{\xi m}^\alpha\right)\dot r_m .
	\label{eq:CN_composition_affinity_projection}
\end{equation}
After the reversible coefficients have been set to zero, the remaining
inequality can be written as
\begin{align}
	0\le\; &
	\sum_m\left(\mc A_m+
	\sum_\xi\sum_\alpha\mu_\xi^\alpha\nu_{\xi m}^\alpha
	-\sum_\xi\sum_\alpha\nu_{\xi m}^\alpha\psi_\xi^\alpha
	\right)\dot r_m
	+\sum_\xi\rho_\xi e_\xi^{\rm int}
		-\sum_\xi\frac{\mbf q_\xi\cdot\nabla_{\mbf x}\theta}{\theta}
		+\mc D_{\rm inelastic} .
	\label{eq:CN_residual_dissipation}
\end{align}
Kinetic crystallization, heat conduction, interphase exchange, plasticity,
remaining terms are nonnegative.


\subsection{Solid mixture rules and nonlinear Biot coefficient}
\label{sec:solid_mixture_rules_nonlinear_biot}

The Coleman--Noll restrictions above close the solid constitutive structure in
terms of one Helmholtz energy for each mineral solid phase.  For \(s\in\mc S\),
the one-component specialization gives
%
% AGENT-LOCK-BEGIN equation-lock: do not edit unless explicitly unlocked by user
\begin{equation}
	\rho_s^s=\rho_s=\phi_s\bar\rho_s^s,
	\qquad
	\eta_s^s=1,
	\qquad
	s\in\mc S .
	\label{eq:solid_final_density_mixture_rules}
\end{equation}
% AGENT-LOCK-END
%
The reference stress transformation is now clean because
\eqref{eq:solid_true_mass_conservation} maps each solid phase intrinsic density
directly through the common true Jacobian.  Define the intrinsic first
Piola--Kirchhoff stress by
%
\begin{equation}
	\bar{\mbf P}_s'
	=
	\bar\rho_{s0}^s
	\left.\p{\psi_s}{\bar{\mbf F}_s^e}\heldfixed{\mc X_s}.
	\label{eq:solid_component_intrinsic_pk_stress_definition}
\end{equation}
%
The single-prime Cauchy stress carried by solid phase \(s\) and its intrinsic
counterpart are
%
% AGENT-LOCK-BEGIN equation-lock: do not edit unless explicitly unlocked by user
\begin{subequations}
	\label{eq:solid_final_stress_mixture_rule}
	\begin{align}
		\sum_{s\in\mc S}\bs\sigma_s'
		 & =\sum_{s\in\mc S}\phi_s\bar{\bs\sigma}_s',
		\label{eq:solid_final_partial_stress_sum_rule}       \\
		\bs\sigma_s'
	 & =
	\rho_s
	\mbf A_s^{-T}
	\left(
		\left.\p{\psi_s}{\bar{\mbf F}_s^e}\heldfixed{\mc X_s}
		(\bar{\mbf F}_s^e)^T
		\right)
		\mbf A_s^T,
		\label{eq:solid_final_component_partial_stress_rule} \\
		\bar{\bs\sigma}_s'
		 & =
		\mbf A_s^{-T}
		\left(
		\bar\rho_s^s
		\left.\p{\psi_s}{\bar{\mbf F}_s^e}\heldfixed{\mc X_s}
		(\bar{\mbf F}_s^e)^T
		\right)
		\mbf A_s^T .
		\label{eq:solid_final_component_stress_rule}
	\end{align}
\end{subequations}
% AGENT-LOCK-END
%
Using \(\bar J_s\bar\rho_s^s=\bar\rho_{s0}^s\), the intrinsic stress has the
standard Piola form inside the distension map:
%
\begin{equation}
	\bar{\bs\sigma}_s'
	=
	\mbf A_s^{-T}
	\left(
	\frac{1}{\bar J_s}
	\bar{\mbf P}_s'
	(\bar{\mbf F}_s^e)^T
	\right)
	\mbf A_s^T .
	\label{eq:solid_modified_piola_transformation}
\end{equation}
%
For the pressure and Biot split below, the Coleman--Noll calculation resolves
solid density dependence through the physical intrinsic density
\(\bar\rho_s^s\) of each solid phase:
%
% AGENT-LOCK-BEGIN equation-lock: do not edit unless explicitly unlocked by user
\begin{empheq}[box=\fbox]{equation}
	\bar p_s=\bar p_s^s
	=\left(\bar\rho_s^s\right)^2
	\left.\p{\psi_s}{\bar\rho_s^s}\heldfixed{\mc X_s},
	\qquad
	s\in\mc S .
	\label{eq:solid_final_pressure_relations}
\end{empheq}
% AGENT-LOCK-END
%
The distension force is the separate Coleman--Noll tensor restriction
%
% AGENT-LOCK-BEGIN equation-lock: do not edit unless explicitly unlocked by user
\begin{equation}
	\sum_{s\in\mc S}\bs\sigma_s'
	=
	\sum_{s\in\mc S}\rho_s
	\left.\p{\psi_s}{\mbf A_s^e}\heldfixed{\mc X_s}
	(\mbf A_s^e)^T .
	\label{eq:solid_final_distension_pressure}
\end{equation}
% AGENT-LOCK-END
%
These are not independent closure equations for double-prime stress.  They are
mixture rules generated by the same Helmholtz energies that supply entropy,
temperature dependence, solid-volume-fraction dependence, and the reversible
distension response.

The same structure gives the overall momentum equation.  Summing the phase
momentum balances, using the mobile-phase pressure relations, and retaining the
conversion-momentum terms gives
%
% AGENT-LOCK-BEGIN equation-lock: do not edit unless explicitly unlocked by user
\begin{equation}
	\begin{aligned}
		\sum_\xi\rho_\xi\mbf a_\xi
		={}& \nabla_{\mbf x}\cdot\left(\sum_{s\in\mc S}\bs\sigma_s'\right)
		+\rho\mbf g+\sum_\xi\mbf b_\xi
		-\sum_{s\in\mc S}\phi_s\nabla_{\mbf x}\bar p_E
		-\sum_l\phi_l\nabla_{\mbf x}\bar p_l \\
		&+\sum_l\sum_\alpha\dot c_l^\alpha\left(\nabla_{\mbf x}\mu_l^\alpha-\mbf v_l\right)
		+\sum_{s\in\mc S}\dot c_s^s\left(\nabla_{\mbf x}\mu_s^s-\mbf v_s\right).
	\end{aligned}
	\label{eq:overall_momentum_phase_pressure_form}
\end{equation}
% AGENT-LOCK-END
%
where \(l\) ranges over the mobile pore phases and \(\rho=\sum_\xi\rho_\xi\).
This is the pressure form of the capillary force.  The
primitive potentials $\gamma_l$ are useful in the variation, but the
volume-fraction Euler--Lagrange relation and the component-storage multiplier
relation combine through \eqref{eq:mobile_capillary_pressure_gradient_collapse}.
The final balance is therefore written in terms of phase pressure gradients,
which is the form used in compositional flow.
For a closed mixture with pairwise internal interactions,
%
% AGENT-LOCK-BEGIN equation-lock: do not edit unless explicitly unlocked by user
\begin{equation}
	\sum_\xi\mbf b_\xi=\mbf 0 .
	\label{eq:internal_momentum_supply_sum}
\end{equation}
% AGENT-LOCK-END
%
The same Euler--Lagrange relations also give the pressure-gradient collapse
needed to compare with the equivalent-pressure form used in multiphase
poromechanics.  Since
$\phi\gamma\left(\left\{S_l\right\}_{l\in\mc M},\theta\right)$ is homogeneous
of degree one in the primitive mobile volume fractions, Euler's theorem gives
$\phi\gamma=\sum_l\phi_l\gamma_l$, or $\gamma=\sum_lS_l\gamma_l$.  Taking the
differential of this Euler relation and subtracting the defining differential
$d(\phi\gamma)=\sum_l\gamma_l\,d\phi_l$ gives the Gibbs--Duhem identity
$\sum_l\phi_l\,d\gamma_l=0$ \cite{callen1985thermodynamics}.  In gradient form,
$\sum_l S_l\nabla_{\mbf x}\gamma_l=\mbf 0$.  Therefore, since
$\bar p_l=\gamma_l+\bar p_E$,
%
% AGENT-LOCK-BEGIN equation-lock: do not edit unless explicitly unlocked by user
\begin{equation}
	\sum_l S_l\nabla_{\mbf x}\bar p_l
	=\nabla_{\mbf x}\bar p_E .
	\label{eq:mobile_pressure_gradient_equivalent_pressure}
\end{equation}
% AGENT-LOCK-END
%
Using this identity in \eqref{eq:overall_momentum_phase_pressure_form} gives
%
% AGENT-LOCK-BEGIN equation-lock: do not edit unless explicitly unlocked by user
\begin{equation}
	\sum_\xi\rho_\xi\mbf a_\xi
	= \nabla_{\mbf x}\cdot\left(\sum_{s\in\mc S}\bs\sigma_s'\right)
	+\rho\mbf g
	-\nabla_{\mbf x}\bar p_E
	+\sum_l\sum_\alpha\dot c_l^\alpha\left(\nabla_{\mbf x}\mu_l^\alpha-\mbf v_l\right)
	+\sum_{s\in\mc S}\dot c_s^s\left(\nabla_{\mbf x}\mu_s^s-\mbf v_s\right).
	\label{eq:overall_momentum_equivalent_pressure_uncollected}
\end{equation}
% AGENT-LOCK-END
%
If external momentum is supplied to the mixture, the omitted term
$\sum_\xi\mbf b_\xi$ is added to the right-hand side.

To connect this result with finite-deformation poromechanics, do the pressure
transform phase-by-phase over the solid minerals.  Unlike a miscible fluid phase,
a mineral solid phase is a physical constituent with its own intrinsic density,
current partial density, and corresponding specific volumes,
%
% AGENT-LOCK-BEGIN equation-lock: do not edit unless explicitly unlocked by user
\begin{subequations}
	\label{eq:solid_intrinsic_specific_volume}
	\begin{align}
		\bar v_s^s
		 & =\frac{1}{\bar\rho_s^s},
		\label{eq:solid_intrinsic_specific_volume_current}   \\
		v_s^s
		 & =\frac{1}{\rho_s},
		\label{eq:solid_component_specific_volume_current}   \\
		v_{s0}^s
		 & =\frac{1}{\rho_{s0}^s},
		\label{eq:solid_component_specific_volume_reference} \\
		\bar v_{s0}^s
		 & =\frac{1}{\bar\rho_{s0}^s}.
		\label{eq:solid_intrinsic_specific_volume_reference}
	\end{align}
\end{subequations}
% AGENT-LOCK-END
%
The Coleman--Noll density coefficient has already identified the thermodynamic
solid pressure in \eqref{eq:solid_final_pressure_relations}.  Equivalently,
since \(\bar v_s^s=1/\bar\rho_s^s\),
%
% AGENT-LOCK-BEGIN equation-lock: do not edit unless explicitly unlocked by user
\begin{equation}
	\bar p_s^s
	=-\left.\p{\psi_s}{\bar v_s^s}\right\vert_{\bar{\mbf F}_s^e,J,\theta}.
	\label{eq:solid_component_pressure_density}
\end{equation}
% AGENT-LOCK-END
%
The storage multiplier may be normalized by the current volume occupied by each
solid phase,
%
% AGENT-LOCK-BEGIN equation-lock: do not edit unless explicitly unlocked by user
\begin{equation}
	\bar p_s^s
	=\frac{\tau_s^s}{\phi_s}.
	\label{eq:solid_component_storage_pressure_definition}
\end{equation}
% AGENT-LOCK-END
%
This is a normalization of the multiplier, not a separate variation with respect
to \(\bar\rho_s^s\).  The density dependence has already been accounted for in
the Coleman--Noll chain rule.  The phase Legendre transform is therefore
%
% AGENT-LOCK-BEGIN equation-lock: do not edit unless explicitly unlocked by user
\begin{subequations}
	\label{eq:solid_legendre_transform_pressure}
	\begin{align}
		\widetilde\psi_s
		 & =\psi_s+\bar p_s^s v_s^s,
		\label{eq:solid_legendre_transform_definition} \\
		\left.\p{\widetilde\psi_s}{\bar p_s^s}\right\vert_{\bar{\mbf F}_s^e,J,\theta}
		 & =v_s^s .
		\label{eq:solid_legendre_transform_pressure_derivative}
	\end{align}
\end{subequations}
% AGENT-LOCK-END
%
The pressure in the volume-fraction constraint appears after the solid phase
pressures are related to the common multiplier.  Since solid phases carry no
mobile capillary potential, \eqref{eq:el_lambda} gives
%
% AGENT-LOCK-BEGIN equation-lock: do not edit unless explicitly unlocked by user
\begin{subequations}
	\label{eq:solid_pressure_equivalent_pressure}
	\begin{align}
		\phi_s\lambda
		 & =-\tau_s^s,
		\label{eq:solid_pressure_lambda_storage_sum} \\
		\bar p_E
		 & \equiv-\lambda
		=\frac{\tau_s^s}{\phi_s}
		=\bar p_s^s,
		\qquad s\in\mc S .
		\label{eq:solid_pressure_equivalent_component_average}
	\end{align}
\end{subequations}
% AGENT-LOCK-END
%
Thus the equivalent pressure is the solid pressure appearing in the
volume-fraction constraint; it is not the pressure used to perform the
phase-level Legendre transforms.  The corresponding stress split follows by
differentiating the phase Legendre transform with respect to the elastic deformation.  The
material component mass relation \eqref{eq:material_mass_2} shows that, when the
conversion and composition variables are held fixed in this derivative, the
density dependence induced by \(\bar{\mbf F}_s^e\) enters through the Jacobian.
Equivalently, the phase specific volume may be treated as
\(v_s^s=v_s^s(J,\bar p_s^s,\theta)\) for this calculation.
Therefore
%
% AGENT-LOCK-BEGIN equation-lock: do not edit unless explicitly unlocked by user
\begin{align}
	\left.\p{v_s^s}{\bar{\mbf F}_s^e}\right\vert_{\bar p_s^s,\theta}
	 & =
	\left.\p{v_s^s}{J}\right\vert_{\bar p_s^s,\theta}
	\p{J}{\bar{\mbf F}_s^e} \\
	 & =
	\left.\p{v_s^s}{J}\right\vert_{\bar p_s^s,\theta}
	J(\bar{\mbf F}_s^e)^{-T}.
	\label{eq:solid_specific_volume_deformation_chain_rule}
\end{align}
% AGENT-LOCK-END
%
At fixed \(\bar p_s^s\), the Legendre transform gives
%
% AGENT-LOCK-BEGIN equation-lock: do not edit unless explicitly unlocked by user
\begin{equation}
	\bar\rho_{s0}^s
	\left.\p{\psi_s}{\bar{\mbf F}_s^e}\heldfixed{\mc X_s}
	=
	\bar\rho_{s0}^s
	\left.\p{\widetilde\psi_s}{\bar{\mbf F}_s^e}\right\vert_{\bar p_s^s,\theta}
	-\bar p_s^s \bar\rho_{s0}^s
	\left.\p{v_s^s}{\bar{\mbf F}_s^e}\right\vert_{\bar p_s^s,\theta}.
	\label{eq:solid_legendre_deformation_derivative}
\end{equation}
% AGENT-LOCK-END
%
The first term is the intrinsic first Piola--Kirchhoff stress defined in
\eqref{eq:solid_component_intrinsic_pk_stress_definition}.  The
Legendre-transformed counterpart is
%
% AGENT-LOCK-BEGIN equation-lock: do not edit unless explicitly unlocked by user
\begin{equation}
	\bar{\mbf P}_s^{\prime\prime}
	=\bar\rho_{s0}^s
	\left.\p{\widetilde\psi_s}{\bar{\mbf F}_s^e}\right\vert_{\bar p_s^s,\theta}.
	\label{eq:solid_component_double_prime_pk_stress_definition}
\end{equation}
% AGENT-LOCK-END
%
Thus
%
% AGENT-LOCK-BEGIN equation-lock: do not edit unless explicitly unlocked by user
\begin{equation}
	\bar{\mbf P}_s'
	=
	\bar{\mbf P}_s^{\prime\prime}
	-\bar p_s^s \bar\rho_{s0}^s
	\left.\p{v_s^s}{J}\right\vert_{\bar p_s^s,\theta}
	J(\bar{\mbf F}_s^e)^{-T}.
	\label{eq:solid_component_pk_stress_split}
\end{equation}
% AGENT-LOCK-END
%
Pushing forward through the true deformation and then through the distension
mapping, with the same $\mbf A^{-T}(\cdot)\mbf A^T$ operation used in
\eqref{eq:solid_final_stress_mixture_rule}, gives
%
% AGENT-LOCK-BEGIN equation-lock: do not edit unless explicitly unlocked by user
\begin{equation}
	\bs\sigma_s'
	=
	\bs\sigma_s^{\prime\prime}
	- \frac{1}{v_{s0}^s}
	\left.\p{v_s^s}{J}\right\vert_{\bar p_s^s,\theta}
	\bar p_s^s\mbf I .
	\label{eq:solid_component_cauchy_stress_push_forward}
\end{equation}
% AGENT-LOCK-END
Here \(\bs\sigma_s^{\prime\prime}\) denotes the current double-prime stress of
solid phase \(s\) after the same distension mapping has been applied.  The
pressure correction remains isotropic because
\(\mbf A^{-T}\mbf I\mbf A^T=\mbf I\).
%
This identifies the solid-phase nonlinear Biot coefficient as
%
% AGENT-LOCK-BEGIN equation-lock: do not edit unless explicitly unlocked by user
\begin{align}
	\Aboxed{B_s^s
	 & \equiv 1 - \frac{1}{v_{s0}^s}
	\left.\p{v_s^s}{J}\right\vert_{\bar p_s^s,\theta}} . \\
	 & \approx 1 - \frac{K^{\star s}}{K^s_s}
	\label{eq:component_nonlinear_biot_coefficient}
\end{align}
% AGENT-LOCK-END
%
With this definition, the transformed component stress split is
%
% AGENT-LOCK-BEGIN equation-lock: do not edit unless explicitly unlocked by user
\begin{equation}
	\bs\sigma_s'
	=\bs\sigma_s^{\prime\prime}
	+\left(1-B_s^s\right)\bar p_s^s\mbf I .
	\label{eq:component_biot_stress_split}
\end{equation}
% AGENT-LOCK-END
%
Substitution into the Coleman--Noll mixture rule gives the aggregate solid stress
split
%
% AGENT-LOCK-BEGIN equation-lock: do not edit unless explicitly unlocked by user
\begin{align}
	\sum_{s\in\mc S}\bs\sigma_s'
	 & =\sum_{s\in\mc S}\bs\sigma_s^{\prime\prime}
	+\sum_{s\in\mc S}
	\left(1-B_s^s\right)\bar p_s^s\mbf I,
	\label{eq:solid_biot_mixture_rule_expanded}     \\
	\sum_{s\in\mc S}\bs\sigma_s'
	 & \equiv \bs\sigma^{\prime\prime}
	+\left(1-B_s\right)\bar p_E\mbf I,
	\label{eq:solid_biot_mixture_stress_definition} \\
	\Aboxed{B_s
	 & =
		1 - \frac{1}{\bar p_E}
		\sum_{s\in\mc S}
		\left(1-B_s^s\right)\bar p_s^s} .
	\label{eq:solid_biot_mixture_coefficient_definition}
\end{align}
% AGENT-LOCK-END
%
The plasticity restrictions in \eqref{eq:CN_residual_dissipation} and
\eqref{eq:CN_generic_plastic_flow_rules} are written in terms of the
single-prime stress.  Therefore, if a constitutive plasticity model is
specified in terms of the double-prime stress, the stress must first be mapped
back through \eqref{eq:solid_biot_mixture_stress_definition}, or through the
phase split \eqref{eq:component_biot_stress_split}, before the plastic
driving stress is evaluated.
%
Thus the solid mixture Biot coefficient is obtained only after the solid phase
pressure contributions have been summed.  In this partial-stress form the phase
stresses already carry the current phase density \(\rho_s\), rather than an
additional component mass-fraction weight.

Combining this equivalent-pressure form with
\eqref{eq:solid_biot_mixture_stress_definition} gives
%
% AGENT-LOCK-BEGIN equation-lock: do not edit unless explicitly unlocked by user
\begin{equation}
	\sum_\xi\rho_\xi\mbf a_\xi
	= \nabla_{\mbf x}\cdot\left[
		\bs\sigma^{\prime\prime}
		+\left(1-B_s\right)
		\bar p_E\mbf I
		\right]
	+\rho\mbf g
	-\nabla_{\mbf x}\bar p_E
	+\sum_l\sum_\alpha\dot c_l^\alpha\left(\nabla_{\mbf x}\mu_l^\alpha-\mbf v_l\right)
	+\sum_{s\in\mc S}\dot c_s^s\left(\nabla_{\mbf x}\mu_s^s-\mbf v_s\right).
	\label{eq:overall_momentum_biot_pressure_uncollected}
\end{equation}
% AGENT-LOCK-END
%
Collecting the equivalent-pressure term into the total stress,
%
% AGENT-LOCK-BEGIN equation-lock: do not edit unless explicitly unlocked by user
\begin{empheq}[box=\fbox]{equation}
	\sum_\xi\rho_\xi\mbf a_\xi
	= \nabla_{\mbf x}\cdot\left[
		\bs\sigma^{\prime\prime}-B_s\bar p_E\mbf I
		\right]
	+\rho\mbf g
	+\sum_l\sum_\alpha\dot c_l^\alpha\left(\nabla_{\mbf x}\mu_l^\alpha-\mbf v_l\right)
	+\sum_{s\in\mc S}\dot c_s^s\left(\nabla_{\mbf x}\mu_s^s-\mbf v_s\right).
	\label{eq:overall_momentum_nonlinear_biot}
\end{empheq}
% AGENT-LOCK-END
%
When inertia and conversion-momentum production are neglected, this has the
standard quasistatic total-stress form, with $\bar p_E$ in place of a
saturation-weighted average pressure.

\subsection{Scalar distension simplifications}
\label{sec:scalar_distension_simplifications}

Drumheller's scalar distension model is recovered here as a specialization of
the tensor theory, not as a consequence of frame invariance.  A scalar
distension can describe only an equal stretch in all directions, possibly with a
rigid reorientation of the pore architecture.  It cannot describe anisotropic
collapse of the pore space.  For example, a circular pore collapsing into an
ellipse requires unequal principal stretches, and therefore requires the full
tensor $\mbf A_s$ or an additional anisotropic structural variable.  The scalar
model is nevertheless useful when the pore-architecture change is assumed
isotropic, as in Drumheller's original reduction~\cite{drumheller2000}.

The distinction from Drumheller's symmetry argument is the placement of the
pressure identification.  In Drumheller's equilibrium reduction, the scalar
distension pressure remains in the equilibrium stress transformation before the
symmetry argument is applied; preserving stress symmetry then restricts the full
distension gradient to a scalar stretch and a rotation.  In the Coleman--Noll
calculation above, the pressure identification
\eqref{eq:CN_solid_pressure_identification} removes the isotropic density-power
term before the elastic-distension coefficient is set to zero.  The remaining
restriction \eqref{eq:solid_final_distension_pressure} is tensorial.  Thus the
scalar reduction below is obtained only by imposing the additional
Drumheller-style specialization, not by the preceding tensor derivation itself.

For that specialization,
%
% AGENT-LOCK-BEGIN equation-lock: do not edit unless explicitly unlocked by user
\begin{equation}
	\mbf A_s^\bullet=(a_s^\bullet)^{1/3}\mbf I,
	\qquad
	\bullet\in\{e,p,\theta\},
	\qquad
	s\in\mc S .
	\label{eq:scalar_distension_factors}
\end{equation}
% AGENT-LOCK-END
%
The scalar distension factors multiply:
%
% AGENT-LOCK-BEGIN equation-lock: do not edit unless explicitly unlocked by user
\begin{equation}
	a_s=a_s^e a_s^p a_s^\theta,
	\qquad
	s\in\mc S .
	\label{eq:scalar_distension_product}
\end{equation}
% AGENT-LOCK-END
%
The true solid deformation becomes
%
% AGENT-LOCK-BEGIN equation-lock: do not edit unless explicitly unlocked by user
\begin{equation}
	\mbf F=a_s^{1/3}\bar{\mbf F}_s,
	\qquad
	s\in\mc S .
	\label{eq:scalar_distension_true_deformation}
\end{equation}
% AGENT-LOCK-END
%
The corresponding Jacobian relation is
%
% AGENT-LOCK-BEGIN equation-lock: do not edit unless explicitly unlocked by user
\begin{equation}
	J=a_s\bar J_s,
	\qquad
	s\in\mc S .
	\label{eq:scalar_distension_jacobian}
\end{equation}
% AGENT-LOCK-END
%
The stress transformation in \eqref{eq:solid_final_stress_mixture_rule}
then loses its tensorial distension mapping:
%
% AGENT-LOCK-BEGIN equation-lock: do not edit unless explicitly unlocked by user
\begin{equation}
	\bs\sigma_s'
	=
	\rho_s
	\left.\p{\psi_s}{\bar{\mbf F}_s^e}\heldfixed{\mc X_s}
	(\bar{\mbf F}_s^e)^T ,
	\qquad
	s\in\mc S .
	\label{eq:scalar_distension_component_stress_rule}
\end{equation}
% AGENT-LOCK-END
%
and we recover the standard definition of the Cauchy stress in terms of the Helmholtz energy.
The mixture rule remains
\(\sum_{s\in\mc S}\bs\sigma_s'\), and the Biot split
\eqref{eq:solid_biot_mixture_stress_definition} is unchanged.  The reason is
that the pressure correction is isotropic, while the scalar distension factors
cancel in the map $\mbf A_s^{-T}(\cdot)\mbf A_s^T$.

The Coleman--Noll distension restriction
\eqref{eq:solid_final_distension_pressure} reduces to the scalar trace
condition
%
% AGENT-LOCK-BEGIN equation-lock: do not edit unless explicitly unlocked by user
\begin{equation}
	\frac{1}{3}\operatorname{tr}\left(\sum_{s\in\mc S}\bs\sigma_s'\right)
	=
	\sum_{s\in\mc S}\rho_s a_s^e
	\left.\p{\psi_s}{a_s^e}\heldfixed{\mc X_s}.
	\label{eq:distension_stress_restriction}
\end{equation}
% AGENT-LOCK-END
%
Using the equivalent-pressure Biot split
\eqref{eq:solid_biot_mixture_stress_definition}, this same condition can be
written as a restriction on the double-prime stress:
%
% AGENT-LOCK-BEGIN equation-lock: do not edit unless explicitly unlocked by user
\begin{equation}
	\frac{1}{3}\operatorname{tr}(\bs\sigma^{\prime\prime})
	=
	\sum_{s\in\mc S}\rho_s a_s^e
	\left.\p{\psi_s}{a_s^e}\heldfixed{\mc X_s}
	-\left(1-B_s\right)\bar p_E .
	\label{eq:double_prime_distension_trace}
\end{equation}
% AGENT-LOCK-END
%
Equations~\eqref{eq:distension_stress_restriction} and
\eqref{eq:double_prime_distension_trace} are not extra balance laws for
$a$; they are the scalar form of the tensor admissibility condition on the
reversible solid stress.

The scalar specialization also reduces the plastic-distension power in
\eqref{eq:CN_residual_dissipation} to a hydrostatic work term:
%
% AGENT-LOCK-BEGIN equation-lock: do not edit unless explicitly unlocked by user
\begin{equation}
	\sum_{s\in\mc S}\bs\sigma_s':
	\mbf A_s^e\dot{\mbf A}_s^p
	(\mbf A_s^p)^{-1}(\mbf A_s^e)^{-1}
	=
	\sum_{s\in\mc S}\frac{1}{3}\operatorname{tr}\left(\bs\sigma_s'\right)
	\frac{\dot a_s^p}{a_s^p}.
	\label{eq:scalar_plastic_distension_power}
\end{equation}
% AGENT-LOCK-END
%
Therefore scalar plastic distension is a compaction or dilation mechanism driven
by the single-prime mean stress.  Any anisotropic pore-collapse mechanism must
return to the tensor residual term in \eqref{eq:CN_residual_dissipation}.

\subsection{Distension stress and crystallization pressure in practice}
\label{sec:distension_stress_crystallization}
\label{sec:constitutive_crystallization_distension}

Crystallization pressure is not introduced here as an additional load in the
momentum equation.  The mechanical response must first be generated by an
admissible solid free energy.  The scalar restrictions in
\eqref{eq:distension_stress_restriction} and
\eqref{eq:double_prime_distension_trace} specify the hydrostatic trace that any
isotropic crystallization-pressure model must reproduce.

For this scalar crystallization-pressure specialization, the volumetric split is
%
% AGENT-LOCK-BEGIN equation-lock: do not edit unless explicitly unlocked by user
\begin{equation}
	\ln J=\ln a_s+\ln\bar J_s,
	\qquad
	s\in\mc S .
	\label{eq:crystallization_volumetric_split}
\end{equation}
% AGENT-LOCK-END
%
Thus a penalty written only in terms of $\ln a_s$ or only in terms of
$\ln(J/a_s)=\ln\bar J_s$ will not, in general, produce the trace required by
Coleman--Noll.  The compatible elastic pressure scale must be tied to
$\ln J$, because $J$ is the product that changes with both the distension
and the deformation Jacobian.

The solid phase Legendre transform was defined in
\eqref{eq:solid_legendre_transform_definition} as
$\widetilde\psi_s=\psi_s+\bar p_s^s v_s^s$.  Thus the
transformed potential generates the double-prime stress, while the
untransformed Helmholtz energy generates the total stress.  The difference
between the two stress measures is exactly the Biot pressure correction in
\eqref{eq:component_biot_stress_split}.  Therefore the Biot pressure contribution
should not be inserted a second time into the transformed volumetric energy.  For
a solid phase \(s\), take the current-volume transformed volumetric stress
potential to be
%
% AGENT-LOCK-BEGIN equation-lock: do not edit unless explicitly unlocked by user
\begin{equation}
	\widetilde W_{\rm vol}^s
	=
	\frac{K^{\star s}}{2}\left(\ln J\right)^2 .
	\label{eq:component_transformed_volumetric_energy}
\end{equation}
% AGENT-LOCK-END
%
The admissibility check is direct.  Holding $\bar J_s$ and $\bar p_s^s$
fixed,
%
% AGENT-LOCK-BEGIN equation-lock: do not edit unless explicitly unlocked by user
\begin{equation}
	\left.
	\p{\widetilde W_{\rm vol}^s}{a_s}
	\right\vert_{\bar J_s,\bar p_s^s}
	=
	\frac{K^{\star s}}{a_s}\ln J .
	\label{eq:elastic_distension_compatibility}
\end{equation}
% AGENT-LOCK-END
%
Therefore
%
% AGENT-LOCK-BEGIN equation-lock: do not edit unless explicitly unlocked by user
\begin{equation}
	a_s
	\left.
	\p{\widetilde W_{\rm vol}^s}{a_s}
	\right\vert_{\bar J_s,\bar p_s^s}
	=
	K^{\star s}\ln J .
	\label{eq:component_transformed_volumetric_trace}
\end{equation}
% AGENT-LOCK-END
%
This is exactly the hydrostatic double-prime stress generated by the transformed
volumetric potential:
%
% AGENT-LOCK-BEGIN equation-lock: do not edit unless explicitly unlocked by user
\begin{equation}
	\frac{1}{3}\operatorname{tr}(\bs\sigma_s^{\prime\prime})
	=
	K^{\star s}\ln J .
	\label{eq:component_double_prime_trace_from_energy}
\end{equation}
% AGENT-LOCK-END
%
To return to the original Helmholtz derivative that appears in the
Coleman--Noll restriction, use the same Legendre/Biot relation in scalar trace
form:
%
% AGENT-LOCK-BEGIN equation-lock: do not edit unless explicitly unlocked by user
\begin{equation}
	\rho_s a_s^e
	\left.\p{\psi_s}{a_s^e}\heldfixed{\mc X_s}
	=
	a_s
	\left.
	\p{\widetilde W_{\rm vol}^s}{a_s}
	\right\vert_{\bar J_s,\bar p_s^s}
	+\left(1-B_s^s\right)\bar p_s^s .
	\label{eq:component_legendre_trace_bridge}
\end{equation}
% AGENT-LOCK-END
%
Substituting \eqref{eq:component_transformed_volumetric_trace} into
\eqref{eq:component_legendre_trace_bridge} gives the component total trace:
%
% AGENT-LOCK-BEGIN equation-lock: do not edit unless explicitly unlocked by user
\begin{equation}
	\frac{1}{3}\operatorname{tr}(\bs\sigma_s')
	=
	K^{\star s}\ln J
	+\left(1-B_s^s\right)\bar p_s^s .
	\label{eq:component_total_trace_biot_compatibility}
\end{equation}
% AGENT-LOCK-END
%
Equivalently, the component Coleman--Noll distension force is
%
% AGENT-LOCK-BEGIN equation-lock: do not edit unless explicitly unlocked by user
\begin{equation}
	\rho_s a_s^e
	\left.\p{\psi_s}{a_s^e}\heldfixed{\mc X_s}
	=
	K^{\star s}\ln J
	+\left(1-B_s^s\right)\bar p_s^s .
	\label{eq:component_cn_trace_satisfied}
\end{equation}
% AGENT-LOCK-END
%
Summing over solid phases gives the total trace condition:
%
% AGENT-LOCK-BEGIN equation-lock: do not edit unless explicitly unlocked by user
\begin{equation}
	\frac{1}{3}\operatorname{tr}\left(\sum_{s\in\mc S}\bs\sigma_s'\right)
	=
	\sum_{s\in\mc S}\rho_s a_s^e
	\left.\p{\psi_s}{a_s^e}\heldfixed{\mc X_s},
	\label{eq:crystallization_total_trace_cn_satisfied}
\end{equation}
% AGENT-LOCK-END
%
which is \eqref{eq:distension_stress_restriction}.  Subtracting the Biot
equivalent-pressure contribution recovers the double-prime condition:
%
% AGENT-LOCK-BEGIN equation-lock: do not edit unless explicitly unlocked by user
\begin{equation}
	\frac{1}{3}\operatorname{tr}(\bs\sigma^{\prime\prime})
	=
	\sum_{s\in\mc S} K^{\star s}\ln J
	=
	\sum_{s\in\mc S}\rho_s a_s^e
	\left.\p{\psi_s}{a_s^e}\heldfixed{\mc X_s}
	-\left(1-B_s\right)\bar p_E,
	\label{eq:crystallization_double_prime_trace_cn_satisfied}
\end{equation}
% AGENT-LOCK-END
%
Thus \eqref{eq:component_transformed_volumetric_energy} conforms to the
Coleman--Noll admissibility condition: the same transformed volumetric potential
that supplies the double-prime stress also supplies the distension force required
by \eqref{eq:distension_stress_restriction}.  Once this stress is used in the
momentum equation, no separate crystallization-pressure force should be added.

There is therefore no second mechanical pressure to introduce.  The volumetric
part of the solid Helmholtz energy is already built into
$\operatorname{tr}(\bs\sigma_s')$, and that stress is the
mechanical quantity that enters the momentum equation.  Crystallization enters
the thermodynamics through the mechanism affinity.  The variational
theory identifies the force conjugate to the mechanism progress rate $\dot r_{(m)}$:
%
% AGENT-LOCK-BEGIN equation-lock: do not edit unless explicitly unlocked by user
\begin{equation}
	\mc{A}_{(m)}=-\sum_\xi\sum_\alpha\mu_\xi^\alpha\nu_{\xi (m)}^\alpha .
	\label{eq:crystallization_affinity_projection_recall}
\end{equation}
% AGENT-LOCK-END
%
The entropy inequality requires the kinetic closure to make this solved driving
force admissible:
%
% AGENT-LOCK-BEGIN equation-lock: do not edit unless explicitly unlocked by user
\begin{equation}
	\mc A_{(m)}\dot r_{(m)}\ge 0 .
	\label{eq:crystallization_admissible_kinetics}
\end{equation}
% AGENT-LOCK-END
%
A simple admissible choice is $\dot r_{(m)}=M_{(m)}\mc A_{(m)}$ with $M_{(m)}\ge0$.  A
one-sided precipitation or dissolution law may be used when the mechanism is not
reversible.

Classical crystallization pressure is therefore best used as a way to prescribe
or diagnose this affinity, not as an added stress.  For a mechanism $(m)$
precipitating solid phase \(s_{(m)}\), define the intrinsic specific volume
%
% AGENT-LOCK-BEGIN equation-lock: do not edit unless explicitly unlocked by user
\begin{equation}
	\bar v_{s_{(m)}}^{s_{(m)}}=\frac{1}{\bar\rho_{s_{(m)}}^{s_{(m)}}} .
	\label{eq:crystallization_intrinsic_specific_volume}
\end{equation}
% AGENT-LOCK-END
%
The equivalent classical pressure is the pressure scale whose excess over the
current component trace stress represents the chemical driving force:
%
% AGENT-LOCK-BEGIN equation-lock: do not edit unless explicitly unlocked by user
\begin{equation}
	\bar p^{s_{(m)}}_c
	=\frac{\mc A_{(m)}}{\bar v_{s_{(m)}}^{s_{(m)}}}
	+\frac{1}{3}\operatorname{tr}(\bs\sigma_{s_{(m)}}') .
	\label{eq:crystallization_affinity_trace_stress}
\end{equation}
% AGENT-LOCK-END
%
Equivalently, a crystallization-pressure law supplies the algebraic affinity
restriction
%
% AGENT-LOCK-BEGIN equation-lock: do not edit unless explicitly unlocked by user
\begin{equation}
	\mc A_{(m)}
	=\bar v_{s_{(m)}}^{s_{(m)}}
	\left(
	\bar p^{s_{(m)}}_c
	-\frac{1}{3}\operatorname{tr}(\bs\sigma_{s_{(m)}}')
	\right).
	\label{eq:crystallization_affinity_volumetric}
\end{equation}
% AGENT-LOCK-END
%
For example, a classical supersaturation law may be written
%
% AGENT-LOCK-BEGIN equation-lock: do not edit unless explicitly unlocked by user
\begin{equation}
	\bar p^{s_{(m)}}_c
	=\frac{R\theta}{V_s^{s_{(m)}}}\ln\Omega_{(m)},
	\label{eq:classical_crystallization_pressure}
\end{equation}
% AGENT-LOCK-END
%
where $V_s^{s_{(m)}}$ is the molar volume of the precipitating solid phase.
Substitution into \eqref{eq:crystallization_affinity_volumetric}
gives the corresponding affinity closure
%
% AGENT-LOCK-BEGIN equation-lock: do not edit unless explicitly unlocked by user
\begin{equation}
	\mc A_{(m)}
	=
	\bar v_{s_{(m)}}^{s_{(m)}}
	\left(
	\frac{R\theta}{V_s^{s_{(m)}}}\ln\Omega_{(m)}
	-\frac{1}{3}\operatorname{tr}(\bs\sigma_{s_{(m)}}')
	\right).
	\label{eq:crystallization_supersaturation_affinity}
\end{equation}
% AGENT-LOCK-END
%
At reversible equilibrium, $\mc A_{(m)}=0$, so the component crystallization
pressure equals the current admissible component trace stress:
%
% AGENT-LOCK-BEGIN equation-lock: do not edit unless explicitly unlocked by user
\begin{equation}
	\bar p^{s_{(m)}}_c
	=
	\frac{1}{3}\operatorname{tr}(\bs\sigma_{s_{(m)}}') .
	\label{eq:distension_equilibrium_crystallization}
\end{equation}
% AGENT-LOCK-END
%
If the affinity and stress are solved instead, the associated saturation ratio
is the diagnostic output
%
% AGENT-LOCK-BEGIN equation-lock: do not edit unless explicitly unlocked by user
\begin{equation}
	\Omega_{(m)}
	=\exp\left(
	\frac{V_s^{s_{(m)}}}{R\theta}
	\left[
		\frac{\mc A_{(m)}}{\bar v_{s_{(m)}}^{s_{(m)}}}
		+\frac{1}{3}\operatorname{tr}(\bs\sigma_{s_{(m)}}')
		\right]
	\right).
	\label{eq:crystallization_saturation_ratio_diagnostic}
\end{equation}
% AGENT-LOCK-END
%
Thus a classical crystallization-pressure formula supplies either an algebraic
equilibrium restriction, a diagnostic output, or a calibration target for the
kinetic law.  The simulator still advances the Helmholtz-generated stresses,
balance laws, conversion constraints, and an admissible law for $\dot r_{(m)}$;
it does not add $\bar p^{\alpha_{(m)}}_c$ as a mechanical load.

\subsection{Model summary, closures, and equation count}
\label{sec:model_summary_closure_count}

This section collects the equations that would be assembled in a simulator.  The
point of the summary is to separate primary unknowns, reversible constitutive
relations generated by the Helmholtz energies, and genuinely dissipative closure
laws.  The component mass balances are
%
% AGENT-LOCK-BEGIN equation-lock: do not edit unless explicitly unlocked by user
\begin{equation}
	\frac{\partial}{\partial t}\left(\phi_\xi\bar\rho_\xi\eta_\xi^\alpha\right)
	+\nabla_{\mbf x}\cdot
	\left(\phi_\xi\bar\rho_\xi\eta_\xi^\alpha\mbf v_\xi\right)
	=\sum_m\nu_{\xi (m)}^\alpha\dot r_{(m)} .
	\label{eq:model_summary_mass_balance}
\end{equation}
% AGENT-LOCK-END
%
For \(s\in\mc S\), this balance is used with \(\eta_s^s=1\), giving
\(\partial_t(\phi_s\bar\rho_s^s)+\nabla_{\mbf x}\cdot(\phi_s\bar\rho_s^s\mbf v_s)
=\sum_m\nu_{s (m)}^s\dot r_{(m)}\).  The phase momentum balances may be kept phasewise in the form
\eqref{eq:el_momentum_uniform}, with the mobile-phase pressure-gradient
specialization shown in \eqref{eq:el_mom_l_simplified}, or summed into
the overall poromechanics form \eqref{eq:overall_momentum_nonlinear_biot}.  The
algebraic volume and mobile-composition constraints are
%
% AGENT-LOCK-BEGIN equation-lock: do not edit unless explicitly unlocked by user
\begin{subequations}
	\label{eq:model_summary_constraints}
	\begin{align}
		\sum_\xi\phi_\xi
		 & =1,
		\label{eq:model_summary_saturation_constraint} \\
		\sum_\alpha\eta_l^\alpha
		 & =1,
		\qquad l\in\mc M, \\
		\eta_s^s
		 & =1,
		\qquad s\in\mc S .
		\label{eq:model_summary_mass_fraction_normalization}
	\end{align}
\end{subequations}
% AGENT-LOCK-END
%
with the corresponding Euler--Lagrange relations
%
% AGENT-LOCK-BEGIN equation-lock: do not edit unless explicitly unlocked by user
\begin{subequations}
	\label{eq:model_summary_pressure_relations}
	\begin{align}
		\lambda+\bar p_l-\gamma_l
		 & =0,
		\qquad l\in\mc M,
		\label{eq:model_summary_mobile_pressure_relation} \\
		\lambda+\bar p_s^s
		 & =0,
		\qquad s\in\mc S,
		\label{eq:model_summary_solid_pressure_relation}  \\
		\sum_\alpha\frac{\tau_l^\alpha}{\phi_l}
		 & =\bar p_l,
		\qquad l\in\mc M, \\
		\frac{\tau_s^s}{\phi_s}
		 & =\bar p_s^s,
		\qquad s\in\mc S .
		\label{eq:model_summary_phase_pressure_relation}
	\end{align}
\end{subequations}
% AGENT-LOCK-END
%
The composition stationarity equations are used in their projected form, with
the redundant normal multiplier $\pi_l$ eliminated explicitly.  For mobile phases
this is \eqref{eq:CN_composition_projected_restriction}.  For solid phases the
composition stationarity is replaced by the specialization
\eqref{eq:CN_solid_composition_projected_restriction}, and solid-volume-fraction
dependence of the stress-free factors is carried by
\eqref{eq:CN_solid_volume_fraction_restriction}.  The intrinsic pressures, solid
stresses, entropy, and distension stress are not independent laws; they are
Coleman--Noll derivatives of the Helmholtz energies summarized by
\eqref{eq:CN_stress_definition},
\eqref{eq:CN_distension_force_definition}, \eqref{eq:CN_phase_pressure_definition},
and \eqref{eq:CN_entropy_definitions}.

The conversion potentials are solved from
\eqref{eq:el_conversion_component}.  After the entropy argument identifies the
uniform source-carried closure
$L_\xi^\alpha=\mu_\xi^\alpha-\psi_\xi$, with \(\psi_s\equiv\psi_s^s\) for
one-component solid phases, the
mechanism affinity is the stoichiometric projection
%
% AGENT-LOCK-BEGIN equation-lock: do not edit unless explicitly unlocked by user
\begin{equation}
	\mc A_{(m)}
	=-\sum_\xi\sum_\alpha\mu_\xi^\alpha\nu_{\xi (m)}^\alpha .
	\label{eq:model_summary_affinity}
\end{equation}
% AGENT-LOCK-END
%
The remaining closure laws are dissipative.  They must be chosen so that
%
% AGENT-LOCK-BEGIN equation-lock: do not edit unless explicitly unlocked by user
\begin{equation}
	\begin{aligned}
			0\le{}\; &
			\sum_m\mc A_{(m)}\dot r_{(m)}
			+\sum_{s\in\mc S}\bs\sigma_s':
			\mbf A_s^e\dot{\mbf A}_s^p
			(\mbf A_s^p)^{-1}(\mbf A_s^e)^{-1}      \\
			         & +\sum_{s\in\mc S}\operatorname{tr}\left[
				\mbf S_s'\dot{\bar{\mbf F}}_s^p
				(\bar{\mbf F}_s^p)^{-1}\right]
		+\sum_{\xi<\zeta}\mbf v_{\xi\zeta}\cdot
		\mbf k_{\xi\zeta}^{-1}\mbf v_{\xi\zeta} \\
		         & +\sum_\xi
		\frac{\nabla_{\mbf x}\theta\cdot\mbf K_\xi\nabla_{\mbf x}\theta}{\theta} .
	\end{aligned}
	\label{eq:model_summary_dissipation}
\end{equation}
% AGENT-LOCK-END
%
For example, one may use kinetic laws
$\dot r_{(m)}=\mc K_{(m)}(\mc A_{(m)},\hbox{state})$ with
$\mc A_{(m)}\dot r_{(m)}\ge0$, plastic flow rules such as
\eqref{eq:CN_generic_plastic_flow_rules}, pairwise interaction laws
\eqref{eq:CN_interphase_exchange_law}, and the heat-flux closure
\eqref{eq:CN_heat_flux_closure}.  A crystallization-pressure law, if used, enters
only by translating it into the affinity through
\eqref{eq:crystallization_affinity_volumetric}.

Let $P_M=|\mc M|$ be the number of mobile phases, $P_S=|\mc S|$ the number of
solid phases, $P=P_M+P_S$ the total number of phases, $N$ the number of mobile
components tracked in each mobile phase, $N_{(m)}$ the number of independent
conversion mechanisms, and $d$ the spatial dimension.  Using phase displacements
as the kinematic unknowns, with velocities obtained by time differentiation, a
convenient fully implicit pointwise algebraic unknown accounting is as follows.
%
This is not a list of state variables only: the mechanism rates $\dot r_{(m)}$ are
process-rate unknowns whose equations are supplied by kinetic laws.
%
\begin{center}
	\begin{tabular}{p{0.72\textwidth}c}
		Unknown family                                                                    & Count     \\
		\hline
		Phase displacements $u_{\xi 1},\ldots,u_{\xi d}$ for $\xi=1,\ldots,P$             & $dP$      \\
		Intrinsic phase densities $\bar\rho_\xi$ for $\xi=1,\ldots,P$                     & $P$       \\
			Mobile mass fractions $\eta_l^1,\ldots,\eta_l^N$ for $l=1,\ldots,P_M$             & $P_MN$    \\
			Phase volume fractions $\phi_\xi$ for $\xi=1,\ldots,P$                            & $P$       \\
			Volume-fraction multiplier $\lambda$                                              & $1$       \\
			Storage multipliers $\tau_l^\alpha$ and $\tau_s^s$                                & $P_MN+P_S$ \\
			Helmholtz chemical potentials $\mu_l^\alpha$ and $\mu_s^s$                        & $P_MN+P_S$ \\
			Mechanism process rates $\dot r_{(1)},\ldots,\dot r_{(N_{(m)})}$                  & $N_{(m)}$ \\
			Solid plastic internal variables $\bar{\mbf F}_s^p$ and $\mbf A_s^p$, when active & $2d^2P_S$ \\
			Temperature $\theta$                                                              & $1$
	\end{tabular}
\end{center}
%
The corresponding scalar unknown count is
%
% AGENT-LOCK-BEGIN equation-lock: do not edit unless explicitly unlocked by user
\begin{equation}
	\begin{aligned}
			N_{\rm unk}
			 & =dP+P+P_MN+P+1+\left(P_MN+P_S\right)
			+\left(P_MN+P_S\right)+N_{(m)}+2d^2P_S+1 \\
			 & =dP+2d^2P_S+3P_MN+2P_S+2P+N_{(m)}+2 ,
	\end{aligned}
	\label{eq:model_summary_unknown_count}
\end{equation}
% AGENT-LOCK-END
%
The same number of scalar equations is supplied by the following blocks:
%
\begin{center}
	\begin{tabular}{p{0.62\textwidth}c}
		Equation block                                                                           & Count     \\
		\hline
		Phase momentum balances \eqref{eq:el_momentum_uniform}                                   & $dP$      \\
		Intrinsic-density relations in \eqref{eq:model_summary_pressure_relations}               & $P$       \\
			Projected mobile composition stationarity \eqref{eq:CN_composition_projected_restriction} & $P_M(N-1)$ \\
			Mobile composition normalizations in \eqref{eq:model_summary_constraints}                & $P_M$     \\
			Volume-fraction Euler--Lagrange relations in \eqref{eq:model_summary_pressure_relations} & $P$       \\
			Volume-fraction constraint in \eqref{eq:model_summary_constraints}                       & $1$       \\
			Component mass balances \eqref{eq:model_summary_mass_balance}                            & $P_MN+P_S$ \\
			Conversion-potential equations \eqref{eq:el_conversion_component}                        & $P_MN+P_S$ \\
			Kinetic laws for $\dot r_{(m)}$                                                          & $N_{(m)}$ \\
			Plastic flow laws for $\bar{\mbf F}_s^p$ and $\mbf A_s^p$, when active                   & $2d^2P_S$ \\
			Energy or temperature equation using \eqref{eq:CN_energy_balance}                        & $1$
	\end{tabular}
\end{center}
%
Thus
%
% AGENT-LOCK-BEGIN equation-lock: do not edit unless explicitly unlocked by user
\begin{equation}
	\begin{aligned}
			N_{\rm eq}
			 & =dP+P+P_M(N-1)+P_M+P+1+\left(P_MN+P_S\right) \\
			 & \quad
			+\left(P_MN+P_S\right)+N_{(m)}+2d^2P_S+1 \\
			 & =dP+2d^2P_S+3P_MN+2P_S+2P+N_{(m)}+2
			=N_{\rm unk}.
	\end{aligned}
	\label{eq:model_summary_equation_count}
\end{equation}
% AGENT-LOCK-END
%
This count treats $\pi_l$, $\mc A_{(m)}$, and $\dot c_\xi^\alpha$ as eliminated
quantities: $\pi_l$ is removed by projecting the mobile composition stationarity
relations onto the $N-1$ dimensional tangent space of
$\sum_\alpha\eta_l^\alpha=1$, $\mc A_{(m)}$ is substituted from
\eqref{eq:model_summary_affinity}, and $\dot c_\xi^\alpha$ is substituted from
the stoichiometric source term in \eqref{eq:model_summary_mass_balance}.  The
solid mass fractions do not add unknowns because \(\eta_s^s=1\) is imposed by the
one-component solid-phase specialization.
